% O Sistema Geodésico Brasileiro — Parte 1: SIRGAS2000. Herança, idealização, campanha e adoção.
% Quatro aulas de 50 minutos (este arquivo, por enquanto, contém a AULA 1; as demais entram nas
% próximas etapas). Mesmo formato, mesma paleta e mesmas convenções de quase_geoide.tex.
% A Parte 2 da série (a rede altimétrica) será outra aula.
%
% Bilíngue (AGENTS.md): o mesmo arquivo gera as duas versões.
%   Português:  latexmk -pdf sgb_sirgas2000.tex      -> sgb_sirgas2000.pdf
%   English:    latexmk -pdf sgb_sirgas2000_en.tex   -> sgb_sirgas2000_en.pdf
%   Limpar auxiliares:  latexmk -c sgb_sirgas2000.tex sgb_sirgas2000_en.tex
% Todo texto visível passa por \tr{português}{english}. Números decimais passam por \num{}, que
% troca vírgula por ponto na versão em inglês. Números derivados: sgb_sirgas2000_numeros.py.
%
% Acompanha o simulador sistemas_coordenadas/ (ECEF, ENU, geodésicas) e continua a aula
% sgr_terrestres.tex (sistema x frame). Não repete o que ela já ensinou.
%
% Fontes dos fatos e números citados (todas lidas na íntegra em 2026-10-13):
%   [A] SIRGAS Final Report, Working Groups I and II (IBGE, jun/2000):
%       geoftp.ibge.gov.br/informacoes_sobre_posicionamento_geodesico/sirgas/report.pdf
%   [B] IBGE, Ajustamento da Rede Planimétrica Brasileira em SIRGAS2000 (jun/2006):
%       .../rede_planialtimetrica/relatorio/rel_sirgas2000.pdf
%   [C] IBGE, Ajustamento da Rede Planimétrica do SGB (SAD 69, 1996/2002):
%       .../rede_planialtimetrica/relatorio/rel_sad69.pdf
%   [D] SIRGAS (sirgas.ipgh.org), páginas "About us", "SIRGAS2000", "WG I", "WG III",
%       "SIRGAS-CON".
%   [E] Decreto n.o 5.334, de 06/01/2005 (Planalto): nova redação do art. 21 do Decreto
%       n.o 89.817/1984.
%   Aula 1 — Córrego Alegre: Hayford 1924, a = 6 378 388 m, 1/f = 297, \xi=\eta=0, N=0, escolhido
%     pelos desvios da vertical medidos em Santa Catarina [C]. PSAD56 (La Canoa, Venezuela, rede
%     HIRAN); 2 113 estações gravimétricas em torno do ponto-datum; Chuá Astro Datum; SAD69:
%     UGGI67, a = 6 378 160 m, 1/f = 298,25, datum Chuá, \xi=0,31", \eta=-3,52", N=0; primeiro
%     ajuste em computador por IAGS/DMA (HAVOC), rede em 10 áreas; anos 80, IBGE com USHER [C].
%     Ajuste global de 1996: 4 760 estações clássicas, 187 GPS, 179 Doppler; BDG em set/1996;
%     variância a posteriori 1,493; deslocamento até 13 m nos extremos [B][C]. Pesos clássicos:
%     direção 3 ppm, distância 1 ppm, azimute astronômico 1 ppm [B]. Rede de 2006: 4 774 clássicas,
%     1 600 GPS, 6 265 estações ajustadas, 16 913 direções, 1 534 distâncias, 389 azimutes [B].
%     Conexões GPS: até 9 m, raio de influência ≈ 500 km; 12 campanhas em 2004-2005 [B]. Marco
%     PREAOCA (1957/1986/1994, 0,796 m) [B]. Levantamentos clássicos cessam no fim dos anos 80;
%     4 receptores GPS no início de 1991 [B]. PMRG criado em 2000 [B].
%   Transformação SAD69 -> SIRGAS2000 (3 translações, 63 estações): ΔX = -67,35 m,
%     ΔY = +3,88 m, ΔZ = -38,22 m [B, tab. 9]. O relatório imprime a latitude astronômica de
%     Chuá como 19°45'14,34"; é erro de digitação (41,34") — o ξ publicado (0,31") confirma.
%   Números derivados (vetor de 78 m, deslocamento por cidade, ppm): sgb_sirgas2000_numeros.py.

\newif\ifEN
\ifdefined\aulaEN \ENtrue \fi

\documentclass[aspectratio=169,11pt]{beamer}

\usetheme[progressbar=frametitle, sectionpage=progressbar, numbering=fraction]{metropolis}
\metroset{block=fill}

\usepackage[T1]{fontenc}
\usepackage{lmodern}
\ifEN
  \usepackage[brazilian,english]{babel}
\else
  \usepackage[english,brazilian]{babel}
\fi
\usepackage{amsmath, amssymb, mathtools}
\usepackage{booktabs}
\usepackage{icomma}
\usepackage{expl3}
\usepackage[most]{tcolorbox}
\usepackage{xspace}
\usepackage{tikz}
\usetikzlibrary{arrows.meta, calc, positioning, decorations.pathreplacing,
                decorations.pathmorphing, patterns, babel}

% ------------------------------------------------------------------ idioma / language
\newcommand{\tr}[2]{\ifEN#2\else#1\fi}
\ExplSyntaxOn
\cs_new_protected:Npn \sgb_format_num:n #1
{
  \tl_set:Nn \l_tmpa_tl {#1}
  \ifEN
    \regex_replace_all:nnN { , } { . } \l_tmpa_tl
  \else
    \regex_replace_all:nnN { \. } { {,} } \l_tmpa_tl
  \fi
  \ensuremath{\l_tmpa_tl}
}
\NewDocumentCommand{\num}{m}{ \sgb_format_num:n {#1} }
\ExplSyntaxOff

% ------------------------------------------------------------------ paleta (a mesma das outras aulas)
\definecolor{mStone}{HTML}{1C1917}
\definecolor{mStoneMid}{HTML}{78716C}
\definecolor{mStoneLight}{HTML}{F5F5F4}
\definecolor{mTeal}{HTML}{0D9488}
\definecolor{mRose}{HTML}{E11D48}
\definecolor{mAmber}{HTML}{D97706}
\definecolor{mIndigo}{HTML}{4F46E5}
\definecolor{mSky}{HTML}{0284C7}

% As ideias que se repetem em toda figura desta aula
\colorlet{corSist}{mIndigo}    % o SISTEMA / idealizado — elipsoide, datum, regra escrita
\colorlet{corFrame}{mTeal}     % o FRAME / medido — marcos, estações, rede
\colorlet{corTempo}{mAmber}    % o TEMPO — época, velocidade (a partir da aula 3)
\colorlet{corErro}{mRose}      % a DIFERENÇA — erro, distorção, deslocamento
\colorlet{corMar}{mSky}        % o GEOIDE — o mar imaginado sob os continentes

\setbeamercolor{normal text}{fg=mStone, bg=white}
\setbeamercolor{palette primary}{fg=white, bg=mStone}
\setbeamercolor{frametitle}{fg=white, bg=mStone}
\setbeamercolor{alerted text}{fg=mRose}
\setbeamercolor{example text}{fg=mTeal}
\setbeamercolor{progress bar}{fg=mTeal, bg=mTeal!20}
\setbeamercolor{progress bar in head/foot}{fg=mTeal, bg=mTeal!20}
\setbeamercolor{progress bar in section page}{fg=mTeal, bg=mTeal!20}
\setbeamercolor{title separator}{fg=mTeal}
\setbeamercolor{block title}{fg=white, bg=mTeal}
\setbeamercolor{block body}{fg=mStone, bg=mTeal!8}

% ------------------------------------------------------------------ caixas
\newtcolorbox{ideia}[1][]{enhanced, colback=mTeal!7, colframe=mTeal, boxrule=0.6pt, arc=2pt,
  left=6pt, right=6pt, top=4pt, bottom=4pt, fonttitle=\bfseries\small, coltitle=white, #1}
\newtcolorbox{cuidado}[1][]{enhanced, colback=mAmber!8, colframe=mAmber, boxrule=0.6pt, arc=2pt,
  left=6pt, right=6pt, top=4pt, bottom=4pt, fonttitle=\bfseries\small, coltitle=white, #1}
\newtcolorbox{armadilha}[1][]{enhanced, colback=mRose!6, colframe=mRose, boxrule=0.6pt, arc=2pt,
  left=6pt, right=6pt, top=4pt, bottom=4pt, fonttitle=\bfseries\small, coltitle=white, #1}
\newtcolorbox{regra}[1][]{enhanced, colback=mIndigo!7, colframe=mIndigo, boxrule=0.6pt, arc=2pt,
  left=6pt, right=6pt, top=4pt, bottom=4pt, fonttitle=\bfseries\small, coltitle=white, #1}

% ------------------------------------------------------------------ estilos TikZ
\tikzset{
  seta/.style={-{Stealth[length=2.4mm]}, thick},
  seta2/.style={{Stealth[length=2mm]}-{Stealth[length=2mm]}},
  rotulo/.style={font=\scriptsize, text=mStone},
  caixa/.style={draw, rounded corners=3pt, thick, align=center, font=\scriptsize},
  elip/.style={corSist, very thick},
  geoide/.style={corMar, very thick},
  rede/.style={corFrame!70, thick},
  marco/.pic={
    \draw[thick, draw=corFrame, fill=corFrame!25] (-0.1,0) -- (0.1,0) -- (0,0.18) -- cycle;
  },
  antena/.pic={
    \draw[thick, draw=corFrame] (0,0) -- (0,0.16);
    \draw[thick, draw=corFrame, fill=corFrame!25] (-0.11,0.16) arc (180:0:0.11) -- cycle;
  },
}

\newcommand{\bnum}[1]{\ifmmode\textcolor{mTeal}{\mathbf{#1}}\else\textbf{\textcolor{mTeal}{#1}}\fi}
\newcommand{\ppm}{\ensuremath{\,\mathrm{ppm}}\xspace}

\title{\tr{O Sistema Geodésico Brasileiro --- Parte 1: SIRGAS2000}{The Brazilian Geodetic
System --- Part 1: SIRGAS2000}}
\subtitle{\tr{Como o Brasil trocou de referencial: a herança, o projeto, a campanha e a adoção
--- quatro aulas}{How Brazil changed its reference system: the inheritance, the design, the
campaign and the adoption --- four classes}}
\author{}
\institute{\tr{Simulador de Sistemas de Coordenadas}{Coordinate Systems Simulator}}
\date{}

\begin{document}

\maketitle

% ==================================================================
\begin{frame}{\tr{O que vamos ver}{What we will see}}
  \setbeamertemplate{section in toc}[sections numbered]
  \small
  \begin{columns}[T]
    \begin{column}{0.56\textwidth}
      \tableofcontents
    \end{column}
    \begin{column}{0.40\textwidth}
      \begin{ideia}[title={\tr{Esta é a Parte 1}{This is Part 1}}]
        \small\tr{A \bnum{posição}: como o Brasil sabe \emph{onde} fica cada ponto, e como mudou
        a forma de saber. A \bnum{altitude} (a rede altimétrica) fica para a Parte
        2.}{\bnum{Position}: how Brazil knows \emph{where} each point is, and how that changed.
        \bnum{Height} (the leveling network) is left for Part 2.}
      \end{ideia}
    \end{column}
  \end{columns}
\end{frame}

\begin{frame}{\tr{Quatro aulas}{Four classes}}
  \begin{columns}[T]
    \begin{column}{0.47\textwidth}
      \begin{regra}[title={\tr{Aula 1 · O problema herdado}{Class 1 · The inherited problem}}]
        \small\tr{Como se media um país inteiro antes do satélite --- e por que o resultado
        envelhecia mal.}{How a whole country was measured before satellites --- and why the
        result aged badly.}
      \end{regra}
      \vspace{2pt}
      \begin{regra}[title={\tr{Aula 3 · A realização}{Class 3 · The realization}}]
        \small\tr{Maio de 2000: 184 estações medidas ao mesmo tempo, e o cálculo que as
        transforma em coordenadas.}{May 2000: 184 stations measured at the same time, and the
        computation that turns them into coordinates.}
      \end{regra}
    \end{column}
    \begin{column}{0.47\textwidth}
      \begin{regra}[title={\tr{Aula 2 · A idealização}{Class 2 · The design}}]
        \small\tr{1993--1995: como se planeja um referencial para um continente, e o ensaio
        geral chamado SIRGAS95.}{1993--1995: how a reference system is planned for a
        continent, and the dress rehearsal called SIRGAS95.}
      \end{regra}
      \vspace{2pt}
      \begin{regra}[title={\tr{Aula 4 · Adoção e vida útil}{Class 4 · Adoption and lifetime}}]
        \small\tr{A lei, a conta de conversão, a rede contínua --- e por que a coordenada
        oficial já não bate com o receptor.}{The law, the conversion formula, the continuous
        network --- and why the official coordinate no longer matches the receiver.}
      \end{regra}
    \end{column}
  \end{columns}
  \medskip
  \centering
  \tr{Cada aula tem 50 minutos, começa relembrando a anterior e termina com um resumo. Vamos
  \bnum{devagar}: uma ideia por slide.}{Each class is 50 minutes long, starts by recalling the
  previous one and ends with a summary. We will go \bnum{slowly}: one idea per slide.}
\end{frame}

% ==================================================================
\begin{frame}[standout]
  \tr{Aula 1}{Class 1}\\[6pt]
  {\normalsize \tr{O problema herdado}{The inherited problem}}\\[4pt]
  {\scriptsize \tr{aula 1 de 4}{class 1 of 4}}
\end{frame}

\section{\tr{Onde paramos}{Where we left off}}

\begin{frame}{\tr{As perguntas de hoje}{Today's questions}}
  \begin{columns}[T]
    \begin{column}{0.58\textwidth}
      \begin{enumerate}
        \setlength{\itemsep}{8pt}
        \item \tr{O que é o Sistema Geodésico Brasileiro, e do que ele é feito?}{What is the
              Brazilian Geodetic System, and what is it made of?}
        \item \tr{Como se media um país inteiro antes do GPS?}{How was a whole country measured
              before GPS?}
        \item \tr{O que é um datum ``topocêntrico'', e por que ele envelhece mal?}{What is a
              ``topocentric'' datum, and why does it age badly?}
        \item \tr{Que erro a rede herdada carregava --- e como se descobriu?}{What error did the
              inherited network carry --- and how was it found?}
      \end{enumerate}
    \end{column}
    \begin{column}{0.38\textwidth}
      \begin{ideia}[title={\tr{Combinado}{Deal}}]
        \tr{No fim da aula, você vai saber responder às quatro --- com as suas
        palavras.}{By the end of the class you will be able to answer all four --- in your own
        words.}
        \medskip

        \tr{Elas voltam no último slide.}{They come back on the last slide.}
      \end{ideia}
    \end{column}
  \end{columns}
\end{frame}

\begin{frame}{\tr{O que a aula anterior deixou}{What the previous class left us}}
  \begin{columns}[T]
    \begin{column}{0.46\textwidth}
      \small
      \tr{Na aula de referenciais terrestres, três ideias:}{In the terrestrial reference
      systems class, three ideas:}
      \begin{enumerate}
        \setlength{\itemsep}{2pt}
        \item \tr{\textcolor{corSist}{\textbf{sistema}} é a regra escrita; não dá para medir
              nada com ele;}{the \textcolor{corSist}{\textbf{system}} is the written rule; you
              cannot measure anything with it;}
        \item \tr{\textcolor{corFrame}{\textbf{frame}} são estações de verdade que obedecem à
              regra com algum errinho;}{the \textcolor{corFrame}{\textbf{frame}} is a set of real
              stations that obey the rule with some small error;}
        \item \tr{uma coordenada precisa dizer \bnum{em qual frame e em qual época}.}{a
              coordinate must say \bnum{in which frame and at which epoch}.}
      \end{enumerate}
      \smallskip
      \tr{Lá vimos o poste que ``anda'' dezenas de metros entre a planta antiga e o mapa de
      hoje. Aqui vamos abrir essa caixa: \alert{de onde vinha aquele deslocamento?}}{There we
      saw the lamp post that ``moves'' tens of meters between the old plan and today's map.
      Here we open that box: \alert{where did that shift come from?}}
    \end{column}
    \begin{column}{0.50\textwidth}
      \centering
      \begin{tikzpicture}[scale=1.0, caixa/.append style={font=\footnotesize}]
        \node[caixa, draw=corSist, fill=corSist!10, text width=2.2cm] (S) at (0,2.2)
          {\tr{sistema}{system}\\[1pt]{\scriptsize\tr{a regra}{the rule}}};
        \node[caixa, draw=corFrame, fill=corFrame!10, text width=2.2cm] (F) at (3.2,2.2)
          {frame\\[1pt]{\scriptsize\tr{as estações}{the stations}}};
        \node[caixa, draw=mStoneMid, fill=white, text width=2.2cm] (C) at (1.6,0.2)
          {\tr{coordenada}{coordinate}\\[1pt]{\scriptsize\tr{frame + época}{frame + epoch}}};
        \draw[seta, mStoneMid] (S) -- node[rotulo, above] {\tr{materializa}{realized by}} (F);
        \draw[seta, mStoneMid] (F) -- (C);
        \node[rotulo, text=corErro, align=center] at (1.6,1.15) {\tr{hoje: o que existia
          \emph{antes} de existir um frame global}{today: what existed \emph{before} there
          was a global frame}};
      \end{tikzpicture}
    \end{column}
  \end{columns}
\end{frame}

\begin{frame}{\tr{O fio da aula}{The thread of the class}}
  \centering
  \begin{tikzpicture}[scale=0.93, caixa/.append style={font=\scriptsize, text width=1.9cm,
      minimum height=1.05cm}]
    \node[caixa, draw=corErro, fill=corErro!14, very thick] (a) at (0,0)
      {\tr{o problema}{the problem}};
    \node[caixa, draw=corSist, fill=corSist!8] (b) at (2.55,0) {\tr{a idealização}{the design}};
    \node[caixa, draw=corSist, fill=corSist!8] (c) at (5.1,0) {\tr{o planejamento}{the plan}};
    \node[caixa, draw=corFrame, fill=corFrame!8] (d) at (7.65,0)
      {\tr{a realização}{the realization}};
    \node[caixa, draw=corFrame, fill=corFrame!8] (e) at (10.2,0) {\tr{a adoção}{the adoption}};
    \node[caixa, draw=corTempo, fill=corTempo!8] (f) at (12.75,0) {\tr{a vida útil}{the
      lifetime}};
    \foreach \x/\y in {a/b, b/c, c/d, d/e, e/f} \draw[seta, mStoneMid] (\x) -- (\y);
    \draw[decorate, decoration={brace, mirror, amplitude=4pt}, mStoneMid]
      (-0.95,-0.75) -- (0.95,-0.75) node[midway, below=4pt, rotulo] {\tr{aula 1}{class 1}};
    \draw[decorate, decoration={brace, mirror, amplitude=4pt}, mStoneMid]
      (1.6,-0.75) -- (6.05,-0.75) node[midway, below=4pt, rotulo] {\tr{aula 2}{class 2}};
    \draw[decorate, decoration={brace, mirror, amplitude=4pt}, mStoneMid]
      (6.7,-0.75) -- (8.6,-0.75) node[midway, below=4pt, rotulo] {\tr{aula 3}{class 3}};
    \draw[decorate, decoration={brace, mirror, amplitude=4pt}, mStoneMid]
      (9.25,-0.75) -- (13.7,-0.75) node[midway, below=4pt, rotulo] {\tr{aula 4}{class 4}};
  \end{tikzpicture}
  \bigskip

  \small
  \tr{Hoje é só o primeiro quadrinho --- mas é o que dá sentido a todos os outros: se o
  problema não ficar claro, o SIRGAS2000 parece um capricho.}{Today is only the first box ---
  but it gives meaning to all the others: if the problem is not clear, SIRGAS2000 looks like a
  whim.}
\end{frame}

% ==================================================================
\section{\tr{O que é o SGB}{What the SGB is}}

\begin{frame}{\tr{O Sistema Geodésico Brasileiro, em uma frase}{The Brazilian Geodetic System,
    in one sentence}}
  \begin{columns}[T]
    \begin{column}{0.46\textwidth}
      \begin{regra}[title={\tr{SGB}{SGB}}]
        \tr{O conjunto de \bnum{estações e redes} que dá ao país suas \bnum{coordenadas},
        \bnum{altitudes} e \bnum{gravidade} oficiais --- estabelecido e mantido pelo
        IBGE.}{The set of \bnum{stations and networks} that gives the country its official
        \bnum{coordinates}, \bnum{heights} and \bnum{gravity} --- established and maintained by
        IBGE.}
      \end{regra}
      \medskip
      \small
      \tr{Não é só uma lista de números: é a \emph{base} sobre a qual se apoiam o mapa, o
      cadastro, a obra e o seu GPS.}{It is not just a list of numbers: it is the \emph{base}
      on which the map, the land register, the construction and your GPS stand.}
      \medskip

      \tr{Em lei: o Decreto 5.334/2005 diz que os referenciais planimétrico e altimétrico da
      cartografia são os do SGB, \emph{conforme estabelecido pelo IBGE}.}{By law: Decree
      5.334/2005 says that the planimetric and altimetric references of the cartography are
      those of the SGB, \emph{as established by IBGE}.}
    \end{column}
    \begin{column}{0.50\textwidth}
      \centering
      \begin{tikzpicture}[scale=0.95, caixa/.append style={font=\scriptsize, text width=2.3cm}]
        \node[caixa, draw=mStone, fill=mStoneLight, minimum width=7.2cm] (top) at (3.4,3.0)
          {\textbf{SGB}};
        \node[caixa, draw=corSist, fill=corSist!10] (p) at (0.6,1.1)
          {\tr{\textbf{planimétrica}}{\textbf{planimetric}}\\ {\tr{onde fica}{where it is}}
           ($\varphi,\lambda$)};
        \node[caixa, draw=corMar, fill=corMar!10] (a) at (3.4,1.1)
          {\tr{\textbf{altimétrica}}{\textbf{altimetric}}\\ \tr{quão alto está}{how high it is}
           ($H$)};
        \node[caixa, draw=mAmber, fill=mAmber!10] (g) at (6.2,1.1)
          {\tr{\textbf{gravimétrica}}{\textbf{gravimetric}}\\ \tr{quanto pesa}{how much it
           weighs} ($g$)};
        \foreach \n in {p,a,g} \draw[thick, mStoneMid] (top.south -| \n) -- (\n.north);
        \node[rotulo, text=corSist, align=center] at (0.6,-0.1)
          {\tr{\textbf{Parte 1}\\ SIRGAS2000}{\textbf{Part 1}\\ SIRGAS2000}};
        \node[rotulo, text=corMar, align=center] at (3.4,-0.1)
          {\tr{\textbf{Parte 2}}{\textbf{Part 2}}};
        \node[rotulo, text=mAmber, align=center] at (6.2,-0.1)
          {\tr{aula de\\ gravimetria}{gravimetry\\ class}};
      \end{tikzpicture}
    \end{column}
  \end{columns}
\end{frame}

\begin{frame}{\tr{Duas famílias de estações: passivas e ativas}{Two families of stations:
    passive and active}}
  \begin{columns}[T]
    \begin{column}{0.47\textwidth}
      \centering
      \begin{tikzpicture}[scale=1.0]
        \foreach \x/\y in {0.2/0.2, 1.3/0.8, 2.3/0.3, 3.2/1.0, 0.7/1.7, 1.9/2.1, 3.0/1.9}
          \pic at (\x,\y) {marco};
        \draw[rede] (0.2,0.3) -- (1.3,0.9) -- (2.3,0.4) -- (3.2,1.1) (0.7,1.8) -- (1.9,2.2)
          -- (3.0,2.0) (0.2,0.3) -- (0.7,1.8) (1.3,0.9) -- (1.9,2.2) (2.3,0.4) -- (3.0,2.0);
        \node[rotulo, text=corFrame, align=center] at (1.7,-0.4)
          {\tr{\textbf{passiva}: marcos no chão\\ que alguém precisa ir ocupar}{\textbf{passive}:
           marks in the ground\\ that someone must go and occupy}};
      \end{tikzpicture}
    \end{column}
    \begin{column}{0.47\textwidth}
      \centering
      \begin{tikzpicture}[scale=1.0]
        \foreach \x/\y in {0.4/0.4, 1.8/1.2, 3.0/0.3, 1.0/2.0, 2.8/1.9} {
          \pic at (\x,\y) {antena};
          \draw[corFrame, thick] (\x+0.16,\y+0.3) arc (-40:40:0.22);
          \draw[corFrame, thick] (\x+0.26,\y+0.34) arc (-40:40:0.32);
        }
        \node[rotulo, text=corFrame, align=center] at (1.7,-0.4)
          {\tr{\textbf{ativa}: receptores ligados\\ 24 horas por dia, todo dia}{\textbf{active}:
           receivers running\\ 24 hours a day, every day}};
      \end{tikzpicture}
    \end{column}
  \end{columns}
  \medskip
  \small
  \tr{A rede \textbf{passiva} é a herança deste capítulo: vértices de triangulação (VT) e
  estações de poligonal (EP). As \textbf{ativas} são a RBMC (receptores GNSS contínuos) e a RMPG
  (marégrafos), que entram na Aula 4 e na Parte 2.}{The \textbf{passive} network is the
  inheritance of this chapter: triangulation vertices (VT) and traverse stations (EP). The
  \textbf{active} ones are the RBMC (continuous GNSS receivers) and the RMPG (tide gauges),
  which come back in Class 4 and in Part 2.}
\end{frame}

\begin{frame}{\tr{Um marco, três placas}{One mark, three plates}}
  \begin{columns}[T]
    \begin{column}{0.46\textwidth}
      \small
      \tr{O vértice \textbf{PREAOCA} foi implantado em \textbf{1957}: uma chapa de metal
      cravada direto na rocha, com a inscrição \texttt{PREAOCA-1957}.}{The vertex
      \textbf{PREAOCA} was set in \textbf{1957}: a metal plate fixed straight into the rock,
      engraved \texttt{PREAOCA-1957}.}
      \medskip

      \tr{Em \textbf{1986} uma nova ocupação encontrou a chapa gravada
      \texttt{PREAOCA-1986}: tinha sido implantada \alert{outra estação}, com o mesmo
      nome.}{In \textbf{1986} a new occupation found the plate engraved
      \texttt{PREAOCA-1986}: \alert{another station} had been set, with the same name.}
      \medskip

      \tr{Em \textbf{1994} o GPS reocupou a placa de 1986 --- e a coordenada obtida diferiu da
      antiga em \bnum{\num{0.796}\,m}.}{In \textbf{1994} GPS reoccupied the 1986 plate --- and
      the coordinate obtained differed from the old one by \bnum{\num{0.796}\,m}.}
    \end{column}
    \begin{column}{0.50\textwidth}
      \centering
      \begin{tikzpicture}[scale=0.95]
        \foreach \x/\y/\t/\c in {0/2.2/{1957}/mStoneMid, 2.6/2.2/{1986}/mAmber,
                                 5.2/2.2/{1994}/corFrame} {
          \draw[thick, fill=mStoneLight, draw=\c, rounded corners=3pt]
            (\x-0.55,\y-0.5) rectangle (\x+0.55,\y+0.5);
          \node[font=\footnotesize\bfseries, text=\c] at (\x,\y+0.15) {\t};
        }
        \node[rotulo] at (0,1.95) {PREAOCA-1957};
        \node[rotulo] at (2.6,1.95) {PREAOCA-1986};
        \node[rotulo] at (5.2,1.95) {PREAOCA-1986};
        \draw[seta, mStoneMid] (0.65,2.2) -- (1.95,2.2);
        \draw[seta, mStoneMid] (3.25,2.2) -- (4.55,2.2);
        \node[rotulo, align=center] at (1.3,2.95) {\tr{estação nova}{new station}};
        \node[rotulo, align=center] at (3.9,2.95) {GPS};
        \draw[corErro, very thick] (0,1.2) -- (5.2,1.2);
        \draw[corErro, very thick] (0,1.1) -- (0,1.3) (5.2,1.1) -- (5.2,1.3);
        \node[rotulo, text=corErro, fill=white, align=center] at (2.6,1.2)
          {\tr{vetor de \num{0.796}\,m}{\num{0.796}\,m vector}};
        \node[rotulo, align=center] at (2.6,0.3)
          {\tr{o IBGE descartou essa conexão do ajuste}{IBGE dropped this connection from the
           adjustment}};
      \end{tikzpicture}
    \end{column}
  \end{columns}
  \medskip
  \begin{ideia}[title={\tr{Lição}{Lesson}}]
    \small\tr{O \textbf{frame} é feito de \bnum{coisas de verdade}: elas são trocadas,
    destruídas e confundidas. (Duas estações brasileiras do SIRGAS foram destruídas antes de
    serem religadas à rede.)}{The \textbf{frame} is made of \bnum{real things}: they get
    replaced, destroyed and mixed up. (Two Brazilian SIRGAS stations were destroyed before being
    tied back to the network.)}
  \end{ideia}
\end{frame}

% ==================================================================
\section{\tr{Medir um país marco a marco}{Measuring a country mark by mark}}

\begin{frame}{\tr{Triangulação: uma régua e muitos ângulos}{Triangulation: one ruler and many
    angles}}
  \begin{columns}[T]
    \begin{column}{0.44\textwidth}
      \small
      \tr{Medir uma distância longa no chão é caro e demorado. Medir \emph{ângulos} é bem
      mais fácil.}{Measuring a long distance on the ground is expensive and slow. Measuring
      \emph{angles} is much easier.}
      \medskip

      \tr{Então: mede-se \bnum{uma só base} com cuidado e, daí em diante, só ângulos. Com a lei
      dos senos, cada triângulo entrega o lado seguinte --- e a cadeia avança.}{So: measure
      \bnum{one base} carefully and from then on only angles. With the law of sines each
      triangle gives the next side --- and the chain moves on.}
      \medskip

      \tr{Para a cadeia não girar, ancora-se a orientação em \textbf{azimutes astronômicos}
      (pelas estrelas).}{To keep the chain from rotating, its orientation is anchored with
      \textbf{astronomical azimuths} (by the stars).}
    \end{column}
    \begin{column}{0.52\textwidth}
      \centering
      \begin{tikzpicture}[scale=0.95]
        \coordinate (P0) at (0,0); \coordinate (P1) at (1.15,1.25); \coordinate (P2) at (2.4,0);
        \coordinate (P3) at (3.55,1.25); \coordinate (P4) at (4.8,0);
        \coordinate (P5) at (5.95,1.25);
        \draw[rede] (P0) -- (P1) -- (P2) -- (P3) -- (P4) -- (P5) (P1) -- (P3) -- (P5)
          (P0) -- (P2) -- (P4);
        \draw[corFrame, line width=2pt] (P0) -- (P2);
        \foreach \p in {P0,P1,P2,P3,P4,P5} \pic at (\p) {marco};
        \node[rotulo, below, text=corFrame] at (1.2,-0.1) {\tr{base medida}{measured base}};
        \draw[corFrame] (P1) ++(-0.28,-0.22) arc (225:325:0.38);
        \node[rotulo, above] at (P1) {\tr{ângulos}{angles}};
        \draw[seta, mAmber] (P5) ++(0.25,0.25) -- ++(0,0.8);
        \node[rotulo, text=mAmber, right] at ($(P5)+(0.1,1.0)$) {\tr{azimute pelas
          estrelas}{azimuth by stars}};
      \end{tikzpicture}
    \end{column}
  \end{columns}
\end{frame}

\begin{frame}{\tr{O que entrava no cálculo}{What went into the computation}}
  \begin{columns}[T]
    \begin{column}{0.54\textwidth}
      \small
      \renewcommand{\arraystretch}{1.35}
      \begin{tabular}{lcr}
        \toprule
        \tr{observação}{observation} & \tr{peso ($\sigma$)}{weight ($\sigma$)} &
          \tr{no ajuste de 2006}{in the 2006 adjustment} \\
        \midrule
        \tr{direção horizontal}{horizontal direction} & \num{3}\ppm & \num{16913} \\
        \tr{distância}{distance} & \num{1}\ppm & \num{1534} \\
        \tr{azimute astronômico}{astronomical azimuth} & \num{1}\ppm & \num{389} \\
        \bottomrule
      \end{tabular}
      \medskip

      \tr{Quase tudo eram \bnum{direções}: ângulos entre alvos. Distâncias eram poucas e
      a orientação vinha de uma lista curta de azimutes.}{Almost everything was
      \bnum{directions}: angles between targets. Distances were few and orientation came from a
      short list of azimuths.}
    \end{column}
    \begin{column}{0.42\textwidth}
      \begin{ideia}[title={\tr{Estações da rede clássica}{Stations of the classical
          network}}]
        \centering
        {\Large\bnum{\num{4774}}}\\[2pt]
        \scriptsize\tr{vértices de triangulação e estações de poligonal, em 2006}{triangulation
        vertices and traverse stations, in 2006}
      \end{ideia}
      \medskip
      \small
      \tr{``\num{3}\ppm'' lê-se \emph{3 partes por milhão}: \num{3}\,mm de erro a cada
      \num{1}\,km.}{``\num{3}\ppm'' reads \emph{3 parts per million}: \num{3}\,mm of error per
      \num{1}\,km.}
    \end{column}
  \end{columns}
\end{frame}

\begin{frame}{\tr{Onde a triangulação não chega}{Where triangulation does not reach}}
  \begin{columns}[T]
    \begin{column}{0.48\textwidth}
      \small
      \tr{Na Amazônia a floresta tapa as linhas de visada: não dá para ver o marco
      seguinte.}{In the Amazon the forest blocks the lines of sight: you cannot see the next
      mark.}
      \medskip

      \tr{A saída, antes do GPS, foram os \bnum{satélites Doppler}. Dois receptores no chão
      escutam o \emph{mesmo} satélite ao mesmo tempo; o que os dois têm em comum se cancela, e
      sobra o \bnum{vetor entre eles}. É a \textbf{translocação}.}{Before GPS the way out was
      the \bnum{Doppler satellites}. Two receivers on the ground listen to the \emph{same}
      satellite at the same time; what they share cancels, and the \bnum{vector between them}
      is left. This is \textbf{translocation}.}
      \medskip

      \tr{A rede de 1996 já trazia \num{179} estações Doppler e \num{15} da rede de
      translocação da Amazônia.}{The 1996 network already had \num{179} Doppler stations and
      \num{15} from the Amazon translocation network.}
    \end{column}
    \begin{column}{0.48\textwidth}
      \centering
      \begin{tikzpicture}[scale=0.95]
        \fill[corFrame!8] (-0.2,-0.3) rectangle (5.4,0.1);
        \node[fill=mAmber!25, draw=mAmber, circle, inner sep=2pt] at (2.7,2.6) {};
        \node[rotulo, above] at (2.7,2.8) {\tr{satélite}{satellite}};
        \pic at (0.9,0.1) {antena};
        \pic at (4.5,0.1) {antena};
        \draw[mStoneMid, dashed] (2.7,2.55) -- (0.9,0.45);
        \draw[mStoneMid, dashed] (2.7,2.55) -- (4.5,0.45);
        \draw[seta, corFrame, line width=1.4pt] (1.15,0.0) -- (4.25,0.0);
        \node[rotulo, text=corFrame, below] at (2.7,-0.05) {\tr{vetor entre as estações}{vector
          between the stations}};
        \draw[pattern=north east lines, pattern color=mTeal!35] (1.5,0.45) rectangle (3.9,0.9);
        \node[rotulo, fill=white] at (2.7,0.67) {\tr{floresta}{forest}};
      \end{tikzpicture}
    \end{column}
  \end{columns}
\end{frame}

\begin{frame}{\tr{O erro se acumula}{The error piles up}}
  \begin{columns}[T]
    \begin{column}{0.44\textwidth}
      \small
      \tr{Cada triângulo herda a posição do anterior. Um errinho de ângulo no começo vira um
      erro maior \emph{depois}, porque o triângulo seguinte parte de um ponto já
      deslocado.}{Each triangle inherits the position of the previous one. A small angle error
      at the start becomes a larger error \emph{later}, because the next triangle starts from an
      already shifted point.}
      \medskip

      \tr{Resultado: o erro \bnum{cresce com a distância} ao ponto de partida.}{Result: the
      error \bnum{grows with the distance} from the starting point.}
      \medskip

      \tr{Num país de \num{4000} km, o ponto de partida fica \emph{muito} longe das
      fronteiras.}{In a country 4\,000\,km across, the starting point is \emph{very} far from
      the borders.}
    \end{column}
    \begin{column}{0.52\textwidth}
      \centering
      \begin{tikzpicture}[scale=0.95]
        \coordinate (Q0) at (0,0.2); \coordinate (Q1) at (1.1,0.7); \coordinate (Q2) at (2.2,0.35);
        \coordinate (Q3) at (3.3,0.9); \coordinate (Q4) at (4.4,0.5); \coordinate (Q5) at (5.5,1.1);
        \draw[rede] (Q0) -- (Q1) -- (Q2) -- (Q3) -- (Q4) -- (Q5);
        \fill[mStone] (-0.09,0.11) rectangle (0.09,0.29);
        \foreach \q/\r in {Q1/0.12, Q2/0.2, Q3/0.3, Q4/0.42, Q5/0.56} {
          \draw[corErro, dashed, thick] (\q) circle (\r);
          \fill[corFrame] (\q) circle (0.045);
        }
        \node[rotulo, align=center] at (0,-0.45) {\tr{ponto de partida}{starting
          point}\\ \tr{(fixo)}{(fixed)}};
        \node[rotulo, text=corErro, align=center] at (5.0,-0.45) {\tr{incerteza
          cresce}{uncertainty grows}};
        \draw[seta, corErro] (3.6,-0.4) -- (4.5,-0.4);
      \end{tikzpicture}
    \end{column}
  \end{columns}
\end{frame}

% ==================================================================
\section{\tr{O datum topocêntrico}{The topocentric datum}}

\begin{frame}{\tr{Pregar o elipsoide na Terra}{Pinning the ellipsoid to the Earth}}
  \begin{columns}[T]
    \begin{column}{0.44\textwidth}
      \small
      \tr{Para calcular na rede é preciso uma superfície lisa e simples: o
      \textcolor{corSist}{\textbf{elipsoide}}. A Terra de verdade, com seus calombos, é
      parecida com o \textcolor{corMar}{\textbf{geoide}}.}{To compute on the network we need a
      smooth, simple surface: the \textcolor{corSist}{\textbf{ellipsoid}}. The real Earth, with
      its lumps, resembles the \textcolor{corMar}{\textbf{geoid}}.}
      \medskip

      \tr{Os dois não coincidem. Então é preciso escolher \alert{onde encostar} um no
      outro.}{The two do not coincide. So one has to choose \alert{where to touch} one against
      the other.}
      \medskip

      \tr{No \bnum{ponto-datum} se impõe: o elipsoide tangencia o geoide ($N=0$, e a vertical
      coincide com a normal: $\xi=\eta=0$). \alert{Longe dali, os dois se afastam.}}{At the
      \bnum{datum point} we impose: the ellipsoid touches the geoid ($N=0$, and the vertical
      coincides with the normal: $\xi=\eta=0$). \alert{Far from it, the two drift apart.}}
    \end{column}
    \begin{column}{0.52\textwidth}
      \centering
      \begin{tikzpicture}[scale=0.95]
        \draw[elip, domain=0:8, samples=60, smooth]
          plot (\x, {2.2 - 0.07*(\x-4)^2});
        \draw[geoide, domain=0:8, samples=80, smooth]
          plot (\x, {2.2 - 0.07*(\x-4)^2 + (\x-4)^2*(0.045 + 0.02*sin(3*\x r))});
        \fill[mStone] (4,2.2) circle (0.07);
        \node[rotulo, below=3pt, align=center] at (4,2.15) {\tr{ponto-datum}{datum point}\\
          $N=0$, $\xi=\eta=0$};
        \node[rotulo, text=corSist, above] at (1.0,1.55) {\tr{elipsoide}{ellipsoid}};
        \node[rotulo, text=corMar, above] at (6.8,2.05) {\tr{geoide}{geoid}};
        \draw[seta2, mStone] (7.0,1.35) -- node[rotulo, right] {$N$} (7.0,1.65);
        \node[rotulo, text=mStoneMid] at (4,0.5) {\tr{esquemático, fora de escala}{schematic,
          not to scale}};
      \end{tikzpicture}
    \end{column}
  \end{columns}
\end{frame}

\begin{frame}{\tr{Córrego Alegre: escolhido pelos desvios da vertical}{Córrego Alegre: chosen
    by the deflections of the vertical}}
  \begin{columns}[T]
    \begin{column}{0.50\textwidth}
      \small
      \tr{Ao implantar a cadeia de triangulação em Santa Catarina, mediram-se
      \textbf{astronomicamente} os desvios da vertical. Eles tendiam para o leste: mais massa a
      oeste, menos a leste. Conclusão: o ponto-datum ficaria melhor no \bnum{planalto}.}{While
      setting the triangulation chain in Santa Catarina, deflections of the vertical were
      measured \textbf{astronomically}. They leaned east: more mass to the west, less to the
      east. Conclusion: the datum point would sit better on the \bnum{plateau}.}
      \medskip

      \tr{Escolheu-se o vértice \textbf{Córrego Alegre} e o elipsoide internacional de Hayford.
      Posição e orientação no ponto foram determinadas \textbf{pelas estrelas}. Os cálculos eram
      feitos com calculadoras mecânicas e tábuas de logaritmos.}{The \textbf{Córrego Alegre}
      vertex and Hayford's international ellipsoid were chosen. Position and orientation at the
      point were determined \textbf{by the stars}. The computations were done with mechanical
      calculators and logarithm tables.}
    \end{column}
    \begin{column}{0.46\textwidth}
      \begin{regra}[title={\tr{Datum Córrego Alegre}{Córrego Alegre datum}}]
        \small
        \tr{elipsoide: \textbf{Hayford 1924}}{ellipsoid: \textbf{Hayford 1924}}\\
        $a=\num{6378388}$\,m\quad $1/f=\num{297}$\\[3pt]
        \tr{ponto-datum: vértice Córrego Alegre}{datum point: Córrego Alegre vertex}\\
        $\varphi=-19^\circ 50' 14{,}91''$\\
        $\xi=\eta=0$,\quad $N=0$\,m
      \end{regra}
    \end{column}
  \end{columns}
\end{frame}

\begin{frame}{\tr{De Córrego Alegre ao SAD69}{From Córrego Alegre to SAD69}}
  \begin{columns}[T]
    \begin{column}{0.52\textwidth}
      \small
      \begin{enumerate}
        \setlength{\itemsep}{4pt}
        \item \tr{Para ter \emph{um} referencial para a América do Sul, com apoio da agência
              cartográfica americana, medem-se astronomicamente vértices de vários países: nasce
              o \textbf{PSAD56}, com origem em La Canoa, na Venezuela. Uma rede de trilateração
              (HIRAN) liga as cadeias do Brasil e da Venezuela.}{To have \emph{one} reference
              for South America, with support from the American mapping agency, vertices in
              several countries are measured astronomically: the \textbf{PSAD56} is born, with
              origin at La Canoa, Venezuela. A trilateration network (HIRAN) links the Brazilian
              and Venezuelan chains.}
        \item \tr{A adaptação ao geoide ali era fraca. Mediram-se \num{2113} estações
              gravimétricas em volta de um novo candidato: o vértice \textbf{Chuá}.}{The fit to
              the geoid there was poor. \num{2113} gravimetric stations were measured around a
              new candidate: the \textbf{Chuá} vertex.}
        \item \tr{Com ele, o \textbf{SAD69} (\emph{South American Datum 1969}).}{With it, the
              \textbf{SAD69} (\emph{South American Datum 1969}).}
      \end{enumerate}
    \end{column}
    \begin{column}{0.44\textwidth}
      \begin{regra}[title={\tr{Datum SAD69}{SAD69 datum}}]
        \small
        \tr{elipsoide: \textbf{UGGI67}}{ellipsoid: \textbf{UGGI67}}\\
        $a=\num{6378160}$\,m\quad $1/f=\num{298.25}$\\[3pt]
        \tr{ponto-datum: vértice Chuá}{datum point: Chuá vertex}\\
        $\xi=\num{0.31}''$,\quad $\eta=-\num{3.52}''$,\quad $N=0$\,m
      \end{regra}
      \medskip
      \begin{ideia}[title={\tr{Repare}{Notice}}]
        \small\tr{Para escolher onde pregar o elipsoide, usou-se \bnum{gravimetria} --- a aula
        de gravimetria volta a aparecer.}{To choose where to pin the ellipsoid,
        \bnum{gravimetry} was used --- the gravimetry class shows up again.}
      \end{ideia}
    \end{column}
  \end{columns}
\end{frame}

\begin{frame}{\tr{Três referenciais lado a lado}{Three references side by side}}
  \centering
  \small
  \renewcommand{\arraystretch}{1.5}
  \begin{tabular}{llll}
    \toprule
     & \textbf{\tr{Córrego Alegre}{Córrego Alegre}} & \textbf{SAD69} & \textbf{SIRGAS2000} \\
    \midrule
    \tr{elipsoide}{ellipsoid} & Hayford 1924 & UGGI67 & GRS80 \\
    $a$ (m) & \num{6378388} & \num{6378160} & \num{6378137} \\
    $1/f$ & \num{297} & \num{298.25} & \num{298.257222101} \\
    \tr{origem}{origin} & \tr{ponto-datum no chão}{datum point on the ground} &
      \tr{ponto-datum no chão}{datum point on the ground} &
      \tr{centro de massa da Terra}{Earth's center of mass} \\
    \tr{como se fixa}{how it is fixed} & \tr{estrelas}{stars} & \tr{estrelas + gravimetria}{stars
      + gravimetry} & \tr{satélites}{satellites} \\
    \bottomrule
  \end{tabular}
  \bigskip

  \tr{As duas primeiras colunas são \alert{topocêntricas}: o ``centro'' é um ponto no chão de
  Minas Gerais. A terceira é \bnum{geocêntrica}.}{The first two columns are
  \alert{topocentric}: the ``center'' is a point on the ground in Minas Gerais. The third is
  \bnum{geocentric}.}
\end{frame}

\begin{frame}{\tr{Topocêntrico = centro fora do lugar}{Topocentric = center out of place}}
  \begin{columns}[T]
    \begin{column}{0.46\textwidth}
      \small
      \tr{Como o SAD69 foi pregado num ponto do chão, o centro do seu elipsoide não coincide
      com o centro de massa da Terra.}{Because the SAD69 was pinned at a point on the ground,
      the center of its ellipsoid does not coincide with the Earth's center of mass.}
      \medskip

      \tr{Pelos parâmetros oficiais (3 translações), a distância entre os dois centros é:}{By
      the official parameters (3 translations), the distance between the two centers is:}
      \[
        \sqrt{67{,}35^2 + 3{,}88^2 + 38{,}22^2} \approx \bnum{\num{77.5}\,m}
      \]
      \tr{Cerca de oitenta metros. E tudo o que se mede \emph{a partir} desse centro herda
      o desvio.}{About eighty meters. And everything measured \emph{from} that center inherits
      the offset.}
    \end{column}
    \begin{column}{0.50\textwidth}
      \centering
      \begin{tikzpicture}[scale=1.0]
        \draw[corFrame, very thick] (0,0) circle (1.65);
        \draw[corSist, very thick, dashed] (-0.45,0.28) circle (1.65);
        \draw[thick, corFrame] (-0.08,-0.08) -- (0.08,0.08) (-0.08,0.08) -- (0.08,-0.08);
        \draw[thick, corSist] (-0.53,0.2) -- (-0.37,0.36) (-0.53,0.36) -- (-0.37,0.2);
        \draw[seta, corErro, line width=1.3pt] (0,0) -- (-0.45,0.28);
        \node[rotulo, text=corFrame, right] at (1.7,0.2) {\tr{SIRGAS2000: geocentro}{SIRGAS2000:
          geocenter}};
        \node[rotulo, text=corSist, left] at (-1.8,1.35) {\tr{SAD69: centro
          deslocado}{SAD69: shifted center}};
        \node[rotulo, text=corErro, below] at (-0.2,-0.12) {$\approx 78$\,m};
        \node[rotulo, text=mStoneMid] at (0,-2.1) {\tr{fora de escala}{not to scale}};
      \end{tikzpicture}
    \end{column}
  \end{columns}
\end{frame}

\begin{frame}{\tr{Pare e pense}{Stop and think}}
  \begin{columns}[T]
    \begin{column}{0.46\textwidth}
      \begin{regra}[title={\tr{Uma pergunta}{A question}}]
        \small
        \tr{Os dois centros estão a \num{77.5}\,m. Mas o poste só ``anda'' entre \num{57} e
        \num{76}\,m no chão do Brasil.}{The two centers are \num{77.5}\,m apart. But the post
        only ``moves'' between \num{57} and \num{76}\,m on the Brazilian ground.}
        \medskip

        \tr{\textbf{Por que o deslocamento horizontal não é sempre \num{77.5}\,m?}}{\textbf{Why
        isn't the horizontal shift always \num{77.5}\,m?}}
      \end{regra}
      \smallskip
      \onslide<2->{
      \begin{ideia}[title={\tr{Resposta}{Answer}}]
        \small
        \tr{O vetor de \num{77.5}\,m tem direção fixa no espaço. Em cada lugar, uma parte dele
        aponta ao longo da \alert{vertical} (vira \emph{altura}) e só o resto fica
        \emph{horizontal}. Como a vertical muda de lugar para lugar, a divisão muda
        também.}{The \num{77.5}\,m vector has a fixed direction in space. At each place, a part of
        it points along the \alert{vertical} (it becomes \emph{height}) and only the rest is
        \emph{horizontal}. Because the vertical changes from place to place, the split changes
        too.}
      \end{ideia}
      }
    \end{column}
    \begin{column}{0.50\textwidth}
      \small
      \centering
      \renewcommand{\arraystretch}{1.3}
      \begin{tabular}{lrrr}
        \toprule
        \tr{lugar}{place} & $\Delta N$ & $\Delta E$ & $\Delta h$ \\
        \midrule
        Fortaleza & $-\num{41.7}$ & $-\num{38.9}$ & $-\num{29.5}$ \\
        Curitiba & $-\num{54.2}$ & $-\num{48.5}$ & $-\num{3.0}$ \\
        \tr{Chuí}{Chuí} & $-\num{55.2}$ & $-\num{51.8}$ & $+\num{8.1}$ \\
        \bottomrule
      \end{tabular}
      \smallskip

      {\scriptsize\tr{metros, SAD69 $\to$ SIRGAS2000, calculado com os 3 parâmetros oficiais.
      \emph{Sempre} para sudoeste.}{meters, SAD69 $\to$ SIRGAS2000, computed with the 3 official
      parameters. \emph{Always} to the southwest.}}
      \smallskip

      \onslide<2->{\tr{Em Fortaleza, boa parte do vetor vira altura; no Sul, quase
      nada.}{In Fortaleza a good part of the vector becomes height; in the south, almost
      none.}}
    \end{column}
  \end{columns}
\end{frame}

% ==================================================================
\section{\tr{O custo da herança}{The cost of the inheritance}}

\begin{frame}{\tr{A rede, costurada em pedaços}{The network, stitched in pieces}}
  \begin{columns}[T]
    \begin{column}{0.50\textwidth}
      \small
      \tr{O primeiro ajuste em computador do SAD69 (agência cartográfica americana, programa
      HAVOC) dividiu a rede brasileira em \bnum{10 áreas}, calculadas em separado: o computador
      da época não aguentava mais.}{The first computer adjustment of the SAD69 (American mapping
      agency, program HAVOC) split the Brazilian network into \bnum{10 areas}, computed
      separately: the computer of the time could not take more.}
      \medskip

      \tr{Nos anos 1980 o IBGE seguiu assim: cada nova área era calculada \emph{prendendo} as
      estações de junção às coordenadas do cálculo anterior.}{In the 1980s IBGE went on like
      this: each new area was computed \emph{holding} the junction stations at the coordinates
      of the previous computation.}
      \medskip

      \tr{Cada emenda guarda o erro da vizinha. O próprio IBGE registra que isso
      \alert{inseriu distorções}.}{Each seam keeps its neighbor's error. IBGE itself records
      that this \alert{introduced distortions}.}
    \end{column}
    \begin{column}{0.46\textwidth}
      \centering
      \begin{tikzpicture}[scale=0.9]
        \foreach \c in {0,1,2,3,4} \foreach \r in {0,1} {
          \pgfmathsetmacro{\k}{ifthenelse(mod(\c+\r,2)==0, 12, 4)}
          \fill[corFrame!\k] (\c*1.25,\r*1.1) rectangle (\c*1.25+1.25,\r*1.1+1.1);
          \draw[corFrame, thick] (\c*1.25,\r*1.1) rectangle (\c*1.25+1.25,\r*1.1+1.1);
        }
        \foreach \x/\y in {1.25/0.35, 1.25/1.45, 2.5/0.7, 2.5/1.8, 3.75/0.4, 3.75/1.55,
                           5.0/0.8, 5.0/1.7, 0.6/1.1, 1.9/1.1, 3.1/1.1, 4.4/1.1, 5.6/1.1}
          \fill[corErro] (\x,\y) circle (0.06);
        \node[rotulo, align=center] at (3.1,2.55) {\tr{10 áreas calculadas uma de cada vez}{10
          areas computed one at a time}};
        \node[rotulo, text=corErro, align=center] at (3.1,-0.4)
          {\tr{$\bullet$ estações de junção: o erro passa de uma área à seguinte}{$\bullet$
           junction stations: the error passes from one area to the next}};
      \end{tikzpicture}
    \end{column}
  \end{columns}
\end{frame}

\begin{frame}{\tr{1996: o primeiro ajuste global}{1996: the first global adjustment}}
  \begin{columns}[T]
    \begin{column}{0.52\textwidth}
      \small
      \tr{Em 1996 o IBGE ajustou \emph{tudo de uma vez}, pela primeira vez. Antes, só havia
      ajustes parciais.}{In 1996 IBGE adjusted \emph{everything at once}, for the first time.
      Before, there were only partial adjustments.}
      \medskip

      \renewcommand{\arraystretch}{1.3}
      \begin{tabular}{lr}
        \toprule
        \tr{estações da rede clássica}{classical network stations} & \num{4760} \\
        \tr{estações GPS}{GPS stations} & \num{187} \\
        \tr{estações Doppler}{Doppler stations} & \num{179} \\
        \bottomrule
      \end{tabular}
      \medskip

      \tr{Resultado publicado em setembro de 1996 no Banco de Dados Geodésicos (BDG), o
      catálogo oficial do SGB.}{Result published in September 1996 in the Geodetic Database
      (BDG), the official catalog of the SGB.}
    \end{column}
    \begin{column}{0.44\textwidth}
      \begin{ideia}[title={\tr{O que o IBGE concluiu}{What IBGE concluded}}]
        \small\tr{Pela primeira vez se obteve ``um retrato consistente da qualidade da
        rede''.}{For the first time a ``consistent picture of the quality of the network'' was
        obtained.}
      \end{ideia}
      \medskip
      \begin{cuidado}[title={\tr{E o que ele viu}{And what it saw}}]
        \small\tr{Comparado ao que estava no BDG, o deslocamento horizontal \bnum{aumenta com
        a distância ao ponto-datum}, chegando a \bnum{\num{13}\,m} nos extremos norte, sul e
        oeste.}{Compared with what was in the BDG, the horizontal shift \bnum{increases with the
        distance from the datum point}, reaching \bnum{\num{13}\,m} at the northern, southern
        and western extremes.}
      \end{cuidado}
    \end{column}
  \end{columns}
\end{frame}

\begin{frame}{\tr{A distorção cresce com a distância}{The distortion grows with distance}}
  \begin{columns}[T]
    \begin{column}{0.46\textwidth}
      \small
      \tr{É a mesma figura do ``erro que se acumula'', agora em escala de país.}{It is the same
      picture as ``the error piles up'', now at country scale.}
      \medskip

      \tr{Perto de Chuá, a rede é boa: ali o elipsoide foi pregado. Longe dali, geoide e
      elipsoide se separam, e o ajuste em pedaços acumulou o resto.}{Near Chuá the network is
      good: that is where the ellipsoid was pinned. Far from it, geoid and ellipsoid separate,
      and the piecewise adjustment piled up the rest.}
    \end{column}
    \begin{column}{0.50\textwidth}
      \centering
      \begin{tikzpicture}[scale=0.95]
        \draw[seta, mStone] (0,0) -- (5.3,0) node[rotulo, below left] {\tr{distância ao ponto-datum}
          {distance from the datum point}};
        \draw[seta, mStone] (0,0) -- (0,3.1) node[rotulo, above right]
          {\tr{deslocamento}{shift}};
        \draw[corErro, very thick, domain=0:4.9, samples=40, smooth]
          plot (\x, {0.1 + 0.115*\x^1.55});
        \fill[corErro] (4.9,{0.1 + 0.115*4.9^1.55}) circle (0.07);
        \node[rotulo, text=corErro, above left, align=center] at (4.9,{0.1 + 0.115*4.9^1.55})
          {\tr{até \num{13}\,m}{up to \num{13}\,m}};
        \fill[mStone] (0,0) circle (0.06);
        \node[rotulo, below] at (0.05,-0.05) {Chuá};
        \node[rotulo, text=mStoneMid] at (3.0,2.8) {\tr{esquemático: só o valor máximo é um dado}
          {schematic: only the maximum is a datum}};
      \end{tikzpicture}
    \end{column}
  \end{columns}
\end{frame}

\begin{frame}{\tr{Sua vez}{Your turn}}
  \begin{columns}[T]
    \begin{column}{0.50\textwidth}
      \begin{regra}[title={\tr{Quantas partes por milhão?}{How many parts per million?}}]
        \small
        \begin{enumerate}
          \item[(a)] \tr{Um extremo do país fica a cerca de \num{2000}\,km de Chuá e acumulou
                \num{13}\,m. Quantos mm por km é isso (em \ppm)?}{One extreme of the country is
                about \num{2000}\,km from Chuá and piled up \num{13}\,m. How many mm per km is
                that (in \ppm)?}
          \item[(b)] \tr{A rede clássica ``prometia'' \num{3}\ppm nas direções. Em
                \num{2000}\,km, quanto seria isso?}{The classical network ``promised''
                \num{3}\ppm on directions. Over \num{2000}\,km, how much would that be?}
        \end{enumerate}
      \end{regra}
    \end{column}
    \begin{column}{0.46\textwidth}
      \onslide<2->{
      \begin{ideia}[title={\tr{Respostas}{Answers}}]
        \small\raggedright
        (a) $\dfrac{13\,\mathrm{m}}{2\,000\,000\,\mathrm{m}} = \bnum{\num{6.5}\ppm}$.
        \medskip

        (b) $\num{3}\times10^{-6}\times\num{2000}\,\mathrm{km} = \bnum{\num{6}\,m}$.
        \medskip

        \tr{O erro observado é da \alert{mesma ordem} do prometido: a rede não era ``ruim'',
        estava no limite do que a técnica permitia numa distância dessas.}{The observed error is
        of the \alert{same order} as the promised one: the network was not ``bad'', it was at the
        limit of what the technique allowed over such a distance.}
      \end{ideia}
      }
    \end{column}
  \end{columns}
\end{frame}

\begin{frame}{\tr{O GPS mostra o que a rede escondia}{GPS shows what the network was hiding}}
  \begin{columns}[T]
    \begin{column}{0.50\textwidth}
      \small
      \tr{Os levantamentos clássicos pararam no fim dos anos 1980. No começo de 1991, o IBGE
      comprou \bnum{quatro receptores GPS}. Em 2006, já eram \bnum{\num{1600}} estações GPS.}{The
      classical surveys stopped in the late 1980s. In early 1991 IBGE bought \bnum{four GPS
      receivers}. By 2006 there were \bnum{\num{1600}} GPS stations.}
      \medskip

      \tr{Em 2004 e 2005 foram feitas \num{12} campanhas GPS \emph{em cima} de marcos da rede
      clássica. Descobriu-se que \bnum{uma só conexão GPS muda as coordenadas dos vizinhos em
      até \num{9}\,m}, num raio de cerca de \bnum{\num{500}\,km}.}{In 2004 and 2005, \num{12} GPS
      campaigns were made \emph{on top of} marks of the classical network. It was found that
      \bnum{a single GPS connection changes the coordinates of its neighbors by up to
      \num{9}\,m}, within a radius of about \bnum{\num{500}\,km}.}
    \end{column}
    \begin{column}{0.46\textwidth}
      \centering
      \begin{tikzpicture}[scale=0.95]
        \draw[corFrame, thick, dashed] (2.4,1.6) circle (1.55);
        \pic at (2.4,1.45) {antena};
        \node[rotulo, text=corFrame, align=center] at (2.4,0.95) {GPS};
        \foreach \a/\d in {20/1.3, 80/1.1, 140/1.35, 200/1.0, 260/1.2, 320/1.4} {
          \fill[mStoneMid] ($(2.4,1.6)+(\a:\d)$) circle (0.05);
          \draw[seta, corErro] ($(2.4,1.6)+(\a:\d)$) -- ++({\a+180}:{0.35});
        }
        \node[rotulo, text=corFrame] at (2.4,3.45) {$\approx 500$\,km};
        \node[rotulo, text=corErro, align=center] at (2.4,-0.25) {\tr{as marcas da rede clássica
          se deslocam}{classical network marks shift}};
      \end{tikzpicture}
    \end{column}
  \end{columns}
  \medskip
  \begin{ideia}[title={\tr{E agora?}{And now?}}]
    \small\tr{Em 2000, o IBGE cria o \textbf{PMRG} (Projeto Mudança do Referencial Geodésico),
    por resolução do Congresso Brasileiro de Cartografia de 1999. A mudança é inevitável. A
    pergunta é: \bnum{para qual referencial?}}{In 2000 IBGE creates the \textbf{PMRG}
    (Geodetic Reference Change Project), by resolution of the 1999 Brazilian Cartography
    Congress. The change is unavoidable. The question is: \bnum{to which
    reference?}}
  \end{ideia}
\end{frame}

% ==================================================================
\section{\tr{Fechamento}{Wrap-up}}

\begin{frame}{\tr{A pergunta que vira projeto}{The question that becomes a project}}
  \begin{columns}[T]
    \begin{column}{0.46\textwidth}
      \Large\alert{\tr{``E se o centro fosse o centro da Terra, medido de uma só vez por
      satélites?''}{``What if the center were the Earth's center, measured all at once by
      satellites?''}}
      \normalsize
      \par\bigskip
      \raggedright
      \tr{Se a distorção vem de pregar o elipsoide num ponto do chão e de emendar pedaços,
      a cura é \bnum{pregá-lo no centro de massa} e \bnum{medir tudo ao mesmo tempo}.}{If the
      distortion comes from pinning the ellipsoid to a point on the ground and splicing pieces,
      the cure is to \bnum{pin it at the center of mass} and \bnum{measure everything at the
      same time}.}
    \end{column}
    \begin{column}{0.50\textwidth}
      \centering
      \begin{tikzpicture}[scale=0.95, caixa/.append style={font=\footnotesize, text width=2.4cm}]
        \node[caixa, draw=corErro, fill=corErro!8] (a) at (0,2.0)
          {\tr{antes}{before}\\[1pt]{\scriptsize\tr{centro no chão\\ pedaços emendados}{center
          on the ground\\ spliced pieces}}};
        \node[caixa, draw=corFrame, fill=corFrame!8] (b) at (3.6,2.0)
          {\tr{depois}{after}\\[1pt]{\scriptsize\tr{centro no geocentro\\ medido de uma
          vez}{center at the geocenter\\ measured at once}}};
        \draw[seta, mStoneMid] (a) -- (b);
        \node[rotulo, text=mStoneMid, align=center] at (1.8,0.55)
          {\tr{Aula 2: quem decidiu isso, e como se planeja}{Class 2: who decided this, and how it
          is planned}};
      \end{tikzpicture}
    \end{column}
  \end{columns}
\end{frame}

\begin{frame}{\tr{Para levar para casa --- aula 1}{Take-home --- class 1}}
  \small
  \begin{enumerate}
    \setlength{\itemsep}{2pt}
    \item \tr{O \bnum{SGB} é o conjunto de estações e redes que dá ao país posição, altitude e
          gravidade oficiais. A rede \textbf{passiva} (marcos no chão) é a herança; as
          \textbf{ativas} (RBMC, RMPG) são o presente.}{The \bnum{SGB} is the set of stations
          and networks that gives the country official position, height and gravity. The
          \textbf{passive} network (marks in the ground) is the inheritance; the \textbf{active}
          ones (RBMC, RMPG) are the present.}
    \item \tr{Media-se o país por \bnum{triangulação}: uma base, muitos ângulos, azimutes pelas
          estrelas --- e satélites Doppler onde a mata impedia. O erro \alert{cresce com a
          distância} ao ponto de partida.}{The country was measured by \bnum{triangulation}: one
          base, many angles, azimuths by the stars --- and Doppler satellites where the forest
          got in the way. The error \alert{grows with distance} from the starting point.}
    \item \tr{Um \bnum{datum topocêntrico} pregava o elipsoide num ponto do chão (Córrego
          Alegre, depois Chuá). Por isso o centro do SAD69 está a cerca de \bnum{\num{78}\,m}
          do centro da Terra.}{A \bnum{topocentric datum} pinned the ellipsoid to a point on the
          ground (Córrego Alegre, later Chuá). So the center of the SAD69 is about
          \bnum{\num{78}\,m} from the Earth's center.}
    \item \tr{A rede foi \bnum{emendada em pedaços} por limitação de computador. O ajuste global de
          1996 mostrou deslocamentos de até \num{13}\,m longe do ponto-datum.}{The network was
          \bnum{spliced from pieces} because of computer limits. The 1996 global adjustment showed
          shifts of up to \num{13}\,m far from the datum point.}
    \item \tr{Conexões GPS em marcos clássicos movem os vizinhos em até \num{9}\,m, num raio de
          \num{500}\,km: \bnum{o GPS não só mede, ele audita}.}{GPS connections on classical
          marks move their neighbors by up to \num{9}\,m, within a \num{500}\,km radius:
          \bnum{GPS does not just measure, it audits}.}
  \end{enumerate}
  \begin{ideia}
    \scriptsize \tr{\textbf{Na próxima aula:} em 1993, em Assunção, um grupo decide pregar o
    referencial no centro da Terra --- e começa a planejar como.}{\textbf{Next class:} in 1993,
    in Asunción, a group decides to pin the reference system at the Earth's center --- and
    begins planning how.}
  \end{ideia}
\end{frame}

% ==================================================================
% AULA 2 — A IDEALIZAÇÃO: O PROJETO SIRGAS, 1993–1995
% ==================================================================
\section{\tr{Aula 2: A idealização (1993--1995)}{Class 2: The Design (1993--1995)}}

\begin{frame}{\tr{O gancho da Aula 2: o continente sem centro comum}{Class 2 Hook: The Continent Without a Common Center}}
  \begin{ideia}
    \tr{\textbf{A pergunta desta aula:} como se planeja um sistema geodésico único para um
    continente inteiro, quando cada país usava um elipsoide diferente e nenhum deles tinha centro
    no centro da Terra?}{\textbf{The question for this class:} how do you plan a single geodetic
    reference system for an entire continent, when each country used a different ellipsoid and none
    of them was centered at the Earth's center?}
  \end{ideia}
  \vspace{0.2cm}
  \begin{itemize}
    \setlength{\itemsep}{6pt}
    \item \tr{Na Aula 1 vimos que o Brasil usava \textbf{Córrego Alegre} e depois \textbf{SAD69},
          amarrados ao planalto mineiro com erros acumulados de até \bnum{\num{13}\,m}.}{In Class 1
          we saw that Brazil used \textbf{Córrego Alegre} and later \textbf{SAD69}, pinned to Minas
          Gerais with accumulated errors reaching \bnum{\num{13}\,m}.}
    \item \tr{Mas a Argentina usava \emph{Campo Inchauspe}; a Venezuela e a Colômbia usavam
          \emph{La Canoa} (PSAD56); o Chile usava \emph{SAD56}.}{Meanwhile Argentina used
          \emph{Campo Inchauspe}; Venezuela and Colombia used \emph{La Canoa} (PSAD56); Chile used
          \emph{SAD56}.}
    \item \tr{Na fronteira entre dois países, uma mesma ponte ou usina hidrelétrica tinha duas
          coordenadas oficiais que divergiam por dezenas ou centenas de metros!}{At the border
          between two nations, the exact same bridge or dam had two official sets of coordinates
          differing by dozens or hundreds of meters!}
  \end{itemize}
\end{frame}

\begin{frame}{\tr{A colcha de retalhos sul-americana em 1993}{The South American Patchwork in 1993}}
  \begin{columns}[onlytextwidth]
    \begin{column}{0.48\textwidth}
      \begin{tcolorbox}[colback=mRose!8, colframe=mRose, title={\tr{Problemas na fronteira}{Border issues}}]
        \scriptsize
        \begin{itemize}
          \setlength{\itemsep}{3pt}
          \item \tr{Marcos de fronteira medidos por dois países vizinhos não fechavam.}{Border
                monuments surveyed by neighbouring countries disagreed.}
          \item \tr{Cada nação tinha seu próprio elipsoide de referência (Hayford 1924, UGGI67,
                Clark 1866).}{Each nation had its own reference ellipsoid (Hayford 1924, UGGI67,
                Clarke 1866).}
          \item \tr{Sistemas de navegação por satélite (GPS) forneciam coordenadas geocêntricas,
                incompatíveis com todas as cartas nacionais.}{Satellite navigation systems (GPS)
                provided geocentric coordinates, conflicting with all national map sheets.}
        \end{itemize}
      \end{tcolorbox}
    \end{column}
    \begin{column}{0.48\textwidth}
      \centering
      \begin{tikzpicture}[scale=0.85]
        % Desenho esquemático de dois elipsoides tangentes a superfícies diferentes
        \draw[thick, corMar] (0,0) arc[start angle=150, end angle=30, radius=3.5] node[right] {\tiny \tr{Geoide}{Geoid}};
        \draw[thick, dashed, corSist] (-0.3,-0.4) arc[start angle=155, end angle=45, radius=3.6]
              node[above, font=\tiny, text=corSist] {\tr{Elipsoide País A}{Ellipsoid Country A}};
        \draw[thick, dotted, mAmber] (0.3,-0.6) arc[start angle=145, end angle=35, radius=3.4]
              node[below, font=\tiny, text=mAmber] {\tr{Elipsoide País B}{Ellipsoid Country B}};
        \fill[corErro] (1.5, 0.82) circle (2.5pt) node[above, font=\scriptsize] {\tr{Marco de fronteira}{Border mark}};
        \draw[seta, corErro] (1.5,0.82) -- (1.1,0.4) node[left, font=\tiny] {\tr{$h_A$}{$h_A$}};
        \draw[seta, corErro] (1.5,0.82) -- (1.9,0.2) node[right, font=\tiny] {\tr{$h_B$}{$h_B$}};
      \end{tikzpicture}\\[4pt]
      \tiny \tr{A mesma feição física projetada em superfícies diferentes com origens distintas.}{The
      same physical feature projected onto different surfaces with separate origins.}
    \end{column}
  \end{columns}
\end{frame}

\begin{frame}{\tr{Assunção, outubro de 1993: a virada histórica}{Asunción, October 1993: The Historic Turning Point}}
  \begin{columns}[onlytextwidth]
    \begin{column}{0.54\textwidth}
      \begin{itemize}
        \setlength{\itemsep}{4pt}
        \scriptsize
        \item \tr{Em outubro de 1993, realizou-se em \textbf{Assunção (Paraguai)} a \emph{Conferência
              Internacional para a Definição de um Referencial Geocêntrico Sul-Americano}.}{In October
              1993, the \emph{International Conference for the Definition of a South American Geocentric
              Datum} took place in \textbf{Asunción (Paraguay)}.}
        \item \tr{Organizada em conjunto por:}{Jointly organized by:}
              \begin{itemize}
                \tiny
                \item \textbf{IAG} (\tr{Associação Internacional de Geodésia}{International Association of Geodesy});
                \item \textbf{IPGH} (\tr{Instituto Pan-Americano de Geografia e História}{Pan-American Institute of Geography and History});
                \item \textbf{DMA} (\emph{Defense Mapping Agency}, \tr{hoje NGA, EUA}{now NGA, USA}).
              \end{itemize}
        \item \tr{Nasce ali o \textbf{SIRGAS}: originalmente \emph{Sistema de Referência Geocêntrico
              para a América do Sul}.}{There \textbf{SIRGAS} was born: originally \emph{South American
              Geocentric Reference System}.}
      \end{itemize}
    \end{column}
    \begin{column}{0.44\textwidth}
      \begin{regra}
        \scriptsize
        \tr{\textbf{A meta definida em Assunção:}\\[3pt]
        Adotar um referencial \bnum{geocêntrico único}, compatível com o ITRS e as técnicas espaciais
        (GPS), servindo de base para todos os países do continente.}{\textbf{The goal set in Asunción:}\\[3pt]
        Adopt a \bnum{single geocentric} reference frame, compatible with ITRS and space techniques
        (GPS), serving as the foundation for all nations on the continent.}
      \end{regra}
      \vspace{0.2cm}
      \begin{cuidado}
        \tiny \tr{\textbf{Evolução do nome:} em 2001, a ONU recomendou estender o sistema às
        Américas. Em nov/2020, mudou para \emph{Sistema de Referência Geodésico para as Américas}.}{
        \textbf{Name evolution:} in 2001, the UN recommended extending the system to the Americas.
        In Nov/2020, it became \emph{Geodetic Reference System for the Americas}.}
      \end{cuidado}
    \end{column}
  \end{columns}
\end{frame}

\begin{frame}{\tr{A estrutura inicial de trabalho (1993)}{The Initial Working Structure (1993)}}
  \begin{columns}[onlytextwidth]
    \begin{column}{0.48\textwidth}
      \begin{tcolorbox}[colback=mTeal!8, colframe=mTeal, title={\tr{Grupo de Trabalho I (GT I)}{Working Group I (WG I)}}]
        \scriptsize
        \textbf{\tr{Sistema de Referência}{Reference System}}\\[3pt]
        \begin{itemize}
          \setlength{\itemsep}{2pt}
          \item \tr{Definir o sistema de referência continental geocêntrico.}{Define the continental geocentric reference system.}
          \item \tr{Planejar, executar e processar a rede continental de alta precisão via GPS.}{Plan, execute and process the high-precision continental GPS network.}
          \item \tr{Materializar o primeiro ``esqueleto'' continental.}{Materialize the initial continental ``backbone''.}
        \end{itemize}
      \end{tcolorbox}
    \end{column}
    \begin{column}{0.48\textwidth}
      \begin{tcolorbox}[colback=mIndigo!8, colframe=mIndigo, title={\tr{Grupo de Trabalho II (GT II)}{Working Group II (WG II)}}]
        \scriptsize
        \textbf{\tr{Datum Geocêntrico}{Geocentric Datum}}\\[3pt]
        \begin{itemize}
          \setlength{\itemsep}{2pt}
          \item \tr{Orientar a integração das redes geodésicas nacionais ao novo sistema.}{Guide the integration of national geodetic networks into the new system.}
          \item \tr{Densificar as estações dentro das fronteiras de cada país membro.}{Densify stations within the borders of each member country.}
          \item \tr{Definir parâmetros de transformação entre os data clássicos e o SIRGAS.}{Define transformation parameters between classical datums and SIRGAS.}
        \end{itemize}
      \end{tcolorbox}
    \end{column}
  \end{columns}
  \vspace{0.25cm}
  \begin{ideia}
    \scriptsize \tr{\textbf{E a altitude?} Em 1993, os GTs cuidavam apenas da geometria tridimensional
    ($X, Y, Z$). A unificação da rede altimétrica física só nasceria em 1997 com o \textbf{GT III}
    (tema central da nossa Parte 2!).}{\textbf{What about height?} In 1993, the WGs handled only
    3D geometry ($X,Y,Z$). The unification of physical heights was created in 1997 with \textbf{WG III}
    (the core focus of our Part 2!).}
  \end{ideia}
\end{frame}

\begin{frame}{\tr{As quatro grandes decisões de projeto}{The Four Great Design Decisions}}
  \centering
  \small \tr{Em vez de improvisar, os geodestas em Assunção tomaram 4 decisões fundamentais:}{Instead
  of improvising, geodesists in Asunción made 4 fundamental decisions:}\\[0.3cm]
  \begin{columns}[onlytextwidth]
    \begin{column}{0.23\textwidth}
      \begin{tcolorbox}[colback=mStoneLight, colframe=mTeal, title=\centering \scriptsize \tr{1. O Sistema}{1. The System}]
        \centering \tiny \textbf{ITRS / IERS}\\[2pt]
        \tr{Padrão científico global aberto}{Global open scientific standard}
      \end{tcolorbox}
    \end{column}
    \begin{column}{0.23\textwidth}
      \begin{tcolorbox}[colback=mStoneLight, colframe=mIndigo, title=\centering \scriptsize \tr{2. O Elipsoide}{2. The Ellipsoid}]
        \centering \tiny \textbf{GRS80}\\[2pt]
        \tr{Parâmetros geocêntricos IUGG 1979}{IUGG 1979 geocentric parameters}
      \end{tcolorbox}
    \end{column}
    \begin{column}{0.23\textwidth}
      \begin{tcolorbox}[colback=mStoneLight, colframe=mAmber, title=\centering \scriptsize \tr{3. A Técnica}{3. The Technique}]
        \centering \tiny \textbf{GPS espacial}\\[2pt]
        \tr{Rastreio contínuo e simultâneo}{Continuous simultaneous tracking}
      \end{tcolorbox}
    \end{column}
    \begin{column}{0.23\textwidth}
      \begin{tcolorbox}[colback=mStoneLight, colframe=mRose, title=\centering \scriptsize \tr{4. A Estrutura}{4. The Structure}]
        \centering \tiny \textbf{Em camadas}\\[2pt]
        \tr{Global $\to$ Continental $\to$ Nacional}{Global $\to$ Continental $\to$ National}
      \end{tcolorbox}
    \end{column}
  \end{columns}
  \vspace{0.3cm}
  \tr{Vamos examinar a lógica por trás de cada uma dessas quatro escolhas.}{Let us inspect the logic
  behind each of these four choices.}
\end{frame}

\begin{frame}{\tr{Decisão 1: Por que o ITRS do IERS?}{Decision 1: Why ITRS from IERS?}}
  \begin{columns}[onlytextwidth]
    \begin{column}{0.52\textwidth}
      \begin{itemize}
        \setlength{\itemsep}{4pt}
        \scriptsize
        \item \textbf{\tr{O que se escolheu:}{What was chosen:}} o \emph{International Terrestrial
              Reference System} (ITRS), mantido pelo IERS.
        \item \textbf{\tr{Alternativas:}{Alternatives:}}
              \begin{itemize}
                \tiny
                \item Criar um sistema sul-americano próprio;
                \item Adotar diretamente o WGS 84 militar.
              \end{itemize}
        \item \textbf{\tr{Por que o ITRS?}{Why ITRS?}}
              \begin{itemize}
                \tiny
                \item É o padrão científico mundial mais rigoroso da astronomia e geodésia.
                \item Não depende de um único país nem de uma única agência militar.
                \item Combina quatro técnicas espaciais independentes: VLBI, SLR, GNSS e DORIS.
              \end{itemize}
      \end{itemize}
    \end{column}
    \begin{column}{0.44\textwidth}
      \begin{ideia}
        \scriptsize
        \tr{\textbf{A definição do ITRS:}\\[2pt]
        -- Origem no centro de massa da Terra (incluindo oceanos e atmosfera);\\[2pt]
        -- Escala métrica no sistema próprio relativístico terrestre;\\[2pt]
        -- Orientação alinhada ao polo e meridiano de Greenwich (BIH 1984.0);\\[2pt]
        -- Evolução temporal sem rotação líquida da crosta (condição \emph{NNR}).}{\textbf{ITRS definition:}\\[2pt]
        -- Origin at Earth's center of mass (including oceans and atmosphere);\\[2pt]
        -- Metric scale in local relativistic terrestrial frame;\\[2pt]
        -- Orientation aligned with BIH 1984.0 pole and meridian;\\[2pt]
        -- Time evolution with no-net-rotation of the crust (\emph{NNR} condition).}
      \end{ideia}
    \end{column}
  \end{columns}
\end{frame}

\begin{frame}{\tr{Pare e pense: por que não usar o WGS 84 diretamente?}{Stop and Think: Why Not Use WGS 84 Directly?}}
  \begin{armadilha}
    \tr{\textbf{Pergunta comum:} ``Os receptores GPS dão coordenadas no WGS 84. Por que a América
    do Sul precisou criar o SIRGAS em vez de simplesmente decretar o WGS 84?''}{\textbf{Common question:}
    ``GPS receivers output WGS 84 coordinates. Why did South America need SIRGAS instead of just
    decreeing WGS 84?''}
  \end{armadilha}
  \vspace{0.2cm}
  \begin{columns}[onlytextwidth]
    \begin{column}{0.48\textwidth}
      \onslide<2->{
      \begin{tcolorbox}[colback=mRose!8, colframe=mRose, title={\tr{WGS 84 (Militar)}{WGS 84 (Military)}}]
        \scriptsize
        \begin{itemize}
          \setlength{\itemsep}{2pt}
          \item \tr{Controlado pelo Departamento de Defesa dos EUA.}{Controlled by the US Department of Defense.}
          \item \tr{Não possui marcos físicos perenes acessíveis a civis na América do Sul.}{Lacks permanent physical monuments accessible to civilians in South America.}
          \item \tr{Suas atualizações são feitas sem consulta prévia às agências civis dos outros países.}{Updates occur without prior consultation with other nations' civil agencies.}
        \end{itemize}
      \end{tcolorbox}
      }
    \end{column}
    \begin{column}{0.48\textwidth}
      \onslide<2->{
      \begin{tcolorbox}[colback=mTeal!8, colframe=mTeal, title={\tr{SIRGAS (Científico / Soberano)}{SIRGAS (Scientific / Sovereign)}}]
        \scriptsize
        \begin{itemize}
          \setlength{\itemsep}{2pt}
          \item \tr{É a realização \textbf{regional civil e aberta} do ITRS nas Américas.}{The \textbf{regional civil and open} realization of ITRS in the Americas.}
          \item \tr{Materializado em marcos e estações contínuas mantidas pelos próprios países (IBGE, etc.).}{Materialized on marks and continuous stations run by national agencies (IBGE, etc.).}
          \item \tr{Dados públicos, auditáveis e consistentes com o IGS.}{Public, auditable data fully consistent with IGS.}
        \end{itemize}
      \end{tcolorbox}
      }
    \end{column}
  \end{columns}
\end{frame}

\begin{frame}{\tr{Decisão 2: O Elipsoide de Referência GRS80}{Decision 2: The GRS80 Reference Ellipsoid}}
  \begin{columns}[onlytextwidth]
    \begin{column}{0.54\textwidth}
      \begin{itemize}
        \setlength{\itemsep}{4pt}
        \scriptsize
        \item \tr{O ITRS é um sistema cartesiano ortogonal tridimensional $(X, Y, Z)$.}{ITRS is a 3D
              orthogonal Cartesian coordinate system $(X, Y, Z)$.}
        \item \tr{Para uso cartográfico e prático (latitude $\varphi$, longitude $\lambda$, altura $h$),
              associa-se ao sistema um \textbf{elipsoide de revolução bi-axial}.}{For mapping and
              engineering ($\varphi, \lambda, h$), a \textbf{biaxial ellipsoid of revolution} is associated.}
        \item \tr{O elipsoide adotado pelo SIRGAS é o \textbf{GRS80} (\emph{Geodetic Reference System 1980}):}{The
              ellipsoid adopted by SIRGAS is \textbf{GRS80} (\emph{Geodetic Reference System 1980}):}
              \begin{itemize}
                \tiny
                \item Semi-eixo maior: $a = \bnum{\num{6378137}}\,\mathrm{m}$;
                \item Achatamento: $1/f = \bnum{\num{298.257222101}}$;
                \item Constante gravitacional geocêntrica: $GM = \num{3986005}\times10^{8}\,\mathrm{m^3/s^2}$;
                \item Velocidade angular: $\omega = \num{7292115}\times10^{-11}\,\mathrm{rad/s}$.
              \end{itemize}
      \end{itemize}
    \end{column}
    \begin{column}{0.42\textwidth}
      \begin{cuidado}
        \scriptsize
        \tr{\textbf{GRS80 versus WGS 84:}\\[3pt]
        Ambos têm exatamente o mesmo semi-eixo maior ($a = \num{6378137}$\,m).\\[3pt]
        Diferem apenas no achatamento por \bnum{\num{0.00000016}}, decorrente da definição do harmônico
        $J_2$.\\[3pt]
        Na superfície da Terra, essa diferença produz no máximo \bnum{\num{0.1}\,mm} na altura elipsoidal
        --- \emph{praticamente desprezível para mapeamento!}}{\textbf{GRS80 versus WGS 84:}\\[3pt]
        Both share the exact same semi-major axis ($a = \num{6378137}$\,m).\\[3pt]
        They differ only in flattening by \bnum{\num{0.00000016}} due to the $J_2$ harmonic definition.\\[3pt]
        On Earth's surface, this shift produces at most \bnum{\num{0.1}\,mm} in ellipsoidal height
        --- \emph{negligible for practical surveying!}}
      \end{cuidado}
    \end{column}
  \end{columns}
\end{frame}

\begin{frame}{\tr{Decisão 3: A Técnica de Medição por Satélite}{Decision 3: The Satellite Surveying Technique}}
  \begin{columns}[onlytextwidth]
    \begin{column}{0.50\textwidth}
      \scriptsize
      \tr{Por que o \textbf{GPS} e não VLBI, SLR ou Doppler?}{Why \textbf{GPS} instead of VLBI, SLR, or Doppler?}
      \begin{itemize}
        \setlength{\itemsep}{3pt}
        \item \textbf{VLBI e SLR} \tr{eram as mais precisas do mundo, mas exigiam radiotelescópios de
              dezenas de toneladas e lasers astronômicos gigantescos; inviável em 50 marcos na floresta.}{were
              the most accurate globally, but required multi-ton radiotelescopes and giant lasers;
              unfeasible at 50 remote marks in the rainforest.}
        \item \textbf{Doppler (TRANSIT)} \tr{já estava defasado nos anos 90, com precisão de metros.}{was
              outdated by the 1990s, offering meter-level precision.}
        \item \textbf{GPS} \tr{oferecia a combinação perfeita:}{offered the ideal combination:}
              \begin{itemize}
                \tiny
                \item Antenas e receptores leves, transportáveis por equipe de campo;
                \item Rastreio simultâneo da fase da onda portadora (L1 e L2);
                \item Precisão relativa milimétrica em linhas de base de milhares de quilômetros.
              \end{itemize}
      \end{itemize}
    \end{column}
    \begin{column}{0.46\textwidth}
      \centering
      \begin{tikzpicture}[scale=0.9]
        % Desenho esquemático de satélites transmitindo para antenas
        \node (sat1) at (-1.5,2.2) [caixa, fill=mAmber!15] {\tiny GPS SV1};
        \node (sat2) at (1.5,2.2) [caixa, fill=mAmber!15] {\tiny GPS SV2};
        \pic (m1) at (-1.5,0) {antena};
        \pic (m2) at (1.5,0) {antena};
        \draw[thick] (-2,0) -- (2,0);
        \node at (-1.5,-0.3) [font=\tiny] {\tr{Estação A}{Station A}};
        \node at (1.5,-0.3) [font=\tiny] {\tr{Estação B}{Station B}};
        \draw[dashed, seta, mTeal] (sat1) -- (-1.5,0.25);
        \draw[dashed, seta, mTeal] (sat1) -- (1.5,0.25);
        \draw[dashed, seta, mIndigo] (sat2) -- (-1.5,0.25);
        \draw[dashed, seta, mIndigo] (sat2) -- (1.5,0.25);
        \draw[seta2, corErro, thick] (-1.3,0) -- (1.3,0) node[midway, below, font=\tiny] {\tr{Vetor de base $\Delta \vec{X}$}{Baseline vector $\Delta \vec{X}$}};
      \end{tikzpicture}\\[3pt]
      \tiny \tr{Fase da portadora + rastreio simultâneo = eliminação de erros comuns.}{Carrier phase
      + simultaneous tracking = cancellation of common errors.}
    \end{column}
  \end{columns}
\end{frame}

\begin{frame}{\tr{Decisão 4: Estrutura em Camadas (A analogia do metrô)}{Decision 4: Layered Hierarchy (The Subway Analogy)}}
  \begin{columns}[onlytextwidth]
    \begin{column}{0.54\textwidth}
      \begin{itemize}
        \setlength{\itemsep}{4pt}
        \scriptsize
        \item \tr{Como conectar um marco de esquina no interior do Paraná ao centro da Terra?}{How
              do you link a street corner monument in rural Brazil to the Earth's center?}
        \item \tr{Não se mede a esquina direto do polo norte! Usa-se uma \textbf{hierarquia em camadas}:}{You
              don't measure the corner mark from the North Pole! You use a \textbf{layered hierarchy}:}
              \begin{enumerate}
                \tiny
                \item \textbf{\tr{Camada 1 (Global):}{Layer 1 (Global):}} Rede IGS mundial (dezenas de
                      estações com VLBI/SLR e órbitas ultra-precisas).
                \item \textbf{\tr{Camada 2 (Continental):}{Layer 2 (Continental):}} Rede SIRGAS
                      (esqueleto de estações de alta estabilidade amarradas ao IGS).
                \item \textbf{\tr{Camada 3 (Nacional):}{Layer 3 (National):}} Redes do SGB / IBGE
                      (RBMC contínua e marcos passivos espalhados pelo território).
                \item \textbf{\tr{Camada 4 (Local):}{Layer 4 (Local):}} O topógrafo no campo amarrado
                      ao marco ou à estação RBMC mais próxima.
              \end{enumerate}
      \end{itemize}
    \end{column}
    \begin{column}{0.42\textwidth}
      \centering
      \begin{tikzpicture}[scale=0.8]
        \draw[thick, fill=corSist!15] (-2.2,2.4) rectangle (2.2,3.1) node[midway, font=\scriptsize, align=center] {\textbf{IGS Global}\\\tiny ITRF / Órbitas precisas};
        \draw[thick, fill=corFrame!20] (-2.0,1.4) rectangle (2.0,2.1) node[midway, font=\scriptsize, align=center] {\textbf{SIRGAS Continental}\\\tiny Esqueleto continental};
        \draw[thick, fill=mAmber!20] (-1.8,0.4) rectangle (1.8,1.1) node[midway, font=\scriptsize, align=center] {\textbf{SGB Nacional (IBGE)}\\\tiny RBMC / Rede Passiva};
        \draw[thick, fill=mRose!15] (-1.6,-0.6) rectangle (1.6,0.1) node[midway, font=\scriptsize, align=center] {\textbf{Topografia Local}\\\tiny Obras, estradas, cadastro};
        \draw[seta] (0,2.4) -- (0,2.1);
        \draw[seta] (0,1.4) -- (0,1.1);
        \draw[seta] (0,0.4) -- (0,0.1);
      \end{tikzpicture}\\[2pt]
      \tiny \tr{Cada camada transfere a exatidão da anterior para uma escala mais densa.}{Each layer
      transfers the accuracy of the higher tier to a denser scale.}
    \end{column}
  \end{columns}
\end{frame}

\begin{frame}{\tr{SIRGAS95: O Primeiro Teste do Continente}{SIRGAS95: The Continent's First Test}}
  \begin{columns}[onlytextwidth]
    \begin{column}{0.50\textwidth}
      \begin{tcolorbox}[colback=mTeal!8, colframe=mTeal, title={\tr{A campanha de maio/junho de 1995}{The May/June 1995 Campaign}}]
        \scriptsize
        \begin{itemize}
          \setlength{\itemsep}{3pt}
          \item \textbf{\tr{Período:}{Period:}} 26 de maio a 4 de junho de 1995 (10 dias consecutivos de 24 horas).
          \item \textbf{\tr{Países:}{Countries:}} 11 países da América do Sul mobilizados.
          \item \textbf{\tr{Estações:}{Stations:}} \bnum{\num{58}} estações principais + \num{9} excêntricas.
          \item \textbf{\tr{Brasil:}{Brazil:}} \bnum{\num{11}} estações fundamentais ocupadas pelo IBGE e parceiros.
        \end{itemize}
      \end{tcolorbox}
    \end{column}
    \begin{column}{0.46\textwidth}
      \begin{itemize}
        \setlength{\itemsep}{4pt}
        \scriptsize
        \item \tr{\textbf{Critérios de escolha das estações:}}{\textbf{Station selection criteria:}}
              \begin{itemize}
                \tiny
                \item Todos os observatórios espaciais existentes na América do Sul (SLR, VLBI, DORIS).
                \item Distribuição geográfica homogênea pelo continente.
                \item Estabilidade geológica comprovada (pilares de concreto sobre rocha sã).
                \item Coincidência com marcos das redes clássicas existentes para calcular parâmetros de conversão.
              \end{itemize}
      \end{itemize}
    \end{column}
  \end{columns}
\end{frame}

\begin{frame}{\tr{As 11 estações brasileiras no SIRGAS95}{The 11 Brazilian Stations in SIRGAS95}}
  \begin{columns}[onlytextwidth]
    \begin{column}{0.54\textwidth}
      \scriptsize
      \tr{Distribuídas estrategicamente para cobrir todo o território:}{Strategically distributed
      to cover the entire national territory:}
      \begin{itemize}
        \setlength{\itemsep}{2pt}
        \tiny
        \item \textbf{Fortaleza (CE):} sítio espacial global (VLBI / SLR / GPS).
        \item \textbf{Brasília (DF):} coração do planalto central.
        \item \textbf{Curitiba (PR):} Universidade Federal do Paraná (UFPR).
        \item \textbf{Presidente Prudente (SP):} UNESP.
        \item \textbf{Viçosa (MG):} Universidade Federal de Viçosa (UFV).
        \item \textbf{Belém (PA) e Manaus (AM):} nós amazônicos.
        \item \textbf{Bom Jesus da Lapa (BA):} interior nordestino.
        \item \textbf{Cuiabá (MT):} centro-oeste.
        \item \textbf{Santa Maria (RS):} extremo sul.
        \item \textbf{Cruzeiro do Sul (AC):} extremo oeste.
      \end{itemize}
    \end{column}
    \begin{column}{0.42\textwidth}
      \begin{ideia}
        \scriptsize
        \tr{\textbf{Um salto na história:}\\[3pt]
        Pela primeira vez, pontos separados por mais de \num{3500}\,km na floresta amazônica e no
        pampa gaúcho foram conectados com incerteza \bnum{abaixo de 1 cm} no mesmo sistema!}{\textbf{A
        historic leap:}\\[3pt]
        For the first time, points over \num{3500}\,km apart in the Amazon and the southern plains
        were linked with uncertainty \bnum{under 1 cm} in the same frame!}
      \end{ideia}
    \end{column}
  \end{columns}
\end{frame}

\begin{frame}{\tr{Desafio pioneiro: a calibração das antenas}{Pioneering Challenge: Antenna Calibration}}
  \begin{columns}[onlytextwidth]
    \begin{column}{0.52\textwidth}
      \begin{itemize}
        \setlength{\itemsep}{4pt}
        \scriptsize
        \item \tr{Em 1995, cada país usava marcas diferentes de receptores e antenas: Ashtech,
              Leica, Trimble e TurboRogue.}{In 1995, each nation used different brands of receivers
              and antennas: Ashtech, Leica, Trimble, and TurboRogue.}
        \item \tr{O sinal do satélite não é captado no topo mecânico da antena, mas no seu
              \textbf{centro de fase eletrônico}.}{The satellite signal is captured not at the antenna's
              physical tip, but at its \textbf{electronic phase center}.}
        \item \tr{O centro de fase varia conforme a elevação e o azimute do satélite (\emph{Phase
              Center Variation --- PCV}). Se misturarmos antenas sem calibrar, o erro pode chegar a
              \bnum{\num{5} a \num{10}\,cm} na vertical!}{The phase center varies with satellite
              elevation and azimuth (\emph{PCV}). Mixing uncalibrated antennas creates errors up to
              \bnum{\num{5} to \num{10}\,cm} vertically!}
      \end{itemize}
    \end{column}
    \begin{column}{0.44\textwidth}
      \begin{regra}
        \scriptsize
        \tr{\textbf{A solução adotada em 1995:}\\[3pt]
        Foram instaladas \bnum{9 estações excêntricas} colocalizadas a poucos metros das antenas
        principais com antenas de tipos diferentes.\\[3pt]
        A medição direta das linhas de base curtas permitiu calibrar os centros de fase relativos
        entre todas as marcas usadas no continente.}{\textbf{The solution adopted in 1995:}\\[3pt]
        \bnum{9 eccentric stations} were set up meters away from main marks using differing antenna
        models.\\[3pt]
        Direct short-baseline surveying calibrated relative phase center offsets across all brands
        used in the campaign.}
      \end{regra}
    \end{column}
  \end{columns}
\end{frame}

\begin{frame}{\tr{O cálculo do SIRGAS95: três centros independentes}{SIRGAS95 Processing: Three Independent Centers}}
  \begin{columns}[onlytextwidth]
    \begin{column}{0.52\textwidth}
      \scriptsize
      \tr{Para assegurar total rigor científico e afastar qualquer suspeita de viés, os dados brutos
      foram processados de forma cega por três centros independentes:}{To ensure full scientific rigor
      and avoid bias, raw data were processed independently by three centers:}
      \begin{enumerate}
        \setlength{\itemsep}{2pt}
        \tiny
        \item \textbf{DGFI} (\emph{Deutsches Geodätisches Forschungsinstitut}, \tr{Alemanha}{Germany}) ---
              \tr{software científico Bernese}{Bernese scientific software};
        \item \textbf{NIMA / DMA} (\emph{National Imagery and Mapping Agency}, EUA) ---
              \tr{software científico GIPSY/OASIS}{GIPSY/OASIS scientific software};
        \item \textbf{IBGE} (\tr{Brasil}{Brazil}) --- \tr{software científico Bernese}{Bernese scientific software}.
      \end{enumerate}
      \vspace{0.1cm}
      \tr{As soluções de rede livre foram combinadas por mínimos quadrados e amarradas ao \textbf{ITRF94}
      na época \bnum{1995,4} através de \bnum{7 estações globais do IGS} com transformação de Helmert.}{
      Free-network solutions were combined by least squares and tied to \textbf{ITRF94} at epoch \bnum{1995.4}
      via \bnum{7 global IGS stations} using a Helmert transformation.}
    \end{column}
    \begin{column}{0.44\textwidth}
      \begin{tcolorbox}[colback=mTeal!8, colframe=mTeal, title={\tr{As 7 âncoras IGS de 1995}{The 7 1995 IGS Anchors}}]
        \tiny
        \begin{itemize}
          \setlength{\itemsep}{2pt}
          \item \textbf{Fortaleza} (\tr{Brasil}{Brazil})
          \item \textbf{Kourou} (\tr{Guiana Francesa}{French Guiana})
          \item \textbf{Arequipa} (\tr{Peru}{Peru})
          \item \textbf{Santiago} (\tr{Chile}{Chile})
          \item \textbf{Easter Island} (\tr{Ilha de Páscoa, Chile}{Easter Island, Chile})
          \item \textbf{Galápagos} (\tr{Equador}{Ecuador})
          \item \textbf{Pamatai} (\tr{Polinésia Francesa}{French Polynesia})
        \end{itemize}
      \end{tcolorbox}
    \end{column}
  \end{columns}
\end{frame}

\begin{frame}{\tr{A surpresa da crosta: a Terra é dinâmica!}{Crustal Surprise: The Dynamic Earth!}}
  \begin{columns}[onlytextwidth]
    \begin{column}{0.50\textwidth}
      \begin{itemize}
        \setlength{\itemsep}{4pt}
        \scriptsize
        \item \tr{A campanha SIRGAS95 terminou em 4 de junho de 1995 com sucesso absoluto: precisão
              milimétrica em todo o continente.}{The SIRGAS95 campaign concluded on June 4, 1995 with
              sub-centimeter accuracy across the continent.}
        \item \tr{Menos de dois meses depois, em \bnum{30 de julho de 1995}, ocorreu o violento
              \textbf{terremoto de Antofagasta} ($M_w = 8{,}0$), no norte do Chile!}{Less than two
              months later, on \bnum{July 30, 1995}, the massive \textbf{Antofagasta earthquake}
              ($M_w = 8.0$) struck northern Chile!}
        \item \tr{Marcos geodésicos foram subitamente deslocados em até \bnum{\num{70}\,cm}! As
              coordenadas calculadas em maio já não batiam mais com a rocha em agosto!}{Geodetic
              monuments suddenly shifted by up to \bnum{\num{70}\,cm}! The coordinates calculated in
              May no longer matched the rock in August!}
      \end{itemize}
    \end{column}
    \begin{column}{0.46\textwidth}
      \begin{armadilha}
        \scriptsize
        \tr{\textbf{A grande lição aprendida:}\\[3pt]
        Um sistema geodésico moderno não pode ser uma lista de coordenadas estáticas fixadas para
        sempre.\\[3pt]
        A placa sul-americana desliza a \bnum{\num{1} a \num{2}\,cm/ano} para noroeste, e as bordas
        andinas deformam continuamente.\\[3pt]
        \textbf{Sem modelo de velocidades, um referencial de alta precisão envelhece rapidamente!}}{
        \textbf{The great lesson learned:}\\[3pt]
        A modern geodetic frame cannot just be a list of static coordinates fixed forever.\\[3pt]
        The South American plate drifts at \bnum{\num{1} to \num{2}\,cm/yr} northwestward, and the
        Andean margin deforms continuously.\\[3pt]
        \textbf{Without velocity models, a high-precision frame ages quickly!}}
      \end{armadilha}
    \end{column}
  \end{columns}
\end{frame}

\begin{frame}{\tr{O que faltava após 1995?}{What Was Missing After 1995?}}
  \begin{columns}[onlytextwidth]
    \begin{column}{0.48\textwidth}
      \begin{tcolorbox}[colback=mTeal!8, colframe=mTeal, title={\tr{O que deu certo em 1995}{What worked in 1995}}]
        \scriptsize
        \begin{itemize}
          \setlength{\itemsep}{2pt}
          \item \tr{Provou a cooperação técnica entre 11 países e agências militares e civis.}{Proved technical cooperation among 11 countries, military and civil agencies.}
          \item \tr{Demonstrou repetibilidade e consistência milimétrica com o ITRF94.}{Demonstrated sub-centimeter consistency with ITRF94.}
          \item \tr{Mostrou que o GPS era a ferramenta definitiva para a geodésia continental.}{Proved GPS was the definitive tool for continental geodesy.}
        \end{itemize}
      \end{tcolorbox}
    \end{column}
    \begin{column}{0.48\textwidth}
      \begin{tcolorbox}[colback=mAmber!8, colframe=mAmber, title={\tr{O que ainda faltava}{What was still missing}}]
        \scriptsize
        \begin{itemize}
          \setlength{\itemsep}{2pt}
          \item \tr{\textbf{Densidade:} \num{58} estações era pouco para um continente de \num{18} milhões de $\mathrm{km}^2$.}{\textbf{Density:} \num{58} stations were too few for an \num{18}M $\mathrm{km}^2$ continent.}
          \item \tr{\textbf{América Central e Caribe:} estavam fora do projeto original.}{\textbf{Central America and Caribbean:} were outside the original scope.}
          \item \tr{\textbf{Altimetria física:} nenhuma ligação formal entre os marégrafos das redes de nivelamento e o SIRGAS.}{\textbf{Physical heights:} no formal tie between tide gauges and SIRGAS.}
        \end{itemize}
      \end{tcolorbox}
    \end{column}
  \end{columns}
  \vspace{0.25cm}
  \begin{ideia}
    \scriptsize \tr{Nasce assim a decisão de realizar uma nova e grandiosa campanha no ano \textbf{2000},
    reocupando as estações antigas e conectando todos os marégrafos e redes da América Latina.}{Thus
    arose the decision to conduct a new, massive campaign in the year \textbf{2000}, reoccupying old
    sites and connecting all tide gauges and networks across Latin America.}
  \end{ideia}
\end{frame}

\begin{frame}{\tr{Sua vez: a conta da deriva tectônica}{Your Turn: Calculating Tectonic Drift}}
  \begin{columns}[T]
    \begin{column}{0.50\textwidth}
      \begin{ideia}[title={\tr{Problema}{Problem}}]
        \scriptsize
        \tr{Uma estação geodésica em Curitiba tem velocidade de placa estimada em
        \bnum{\num{1.2}\,cm/ano} para noroeste.}{A geodetic station in Curitiba has an estimated plate motion
        of \bnum{\num{1.2}\,cm/yr} northwestward.}
        \begin{enumerate}
          \setlength{\itemsep}{2pt}
          \item[(a)] \tr{Se medirmos essa estação hoje, quantos centímetros ela se moveu em relação à sua
                coordenada da época 1995,4 após \num{10} anos?}{If we measure this mark today, how many centimeters has it shifted relative to its epoch
                1995.4 coordinate after \num{10} years?}
          \item[(b)] \tr{Se ignorarmos esse movimento num levantamento de precisão em 2005, o erro é
                significativo?}{If we ignore this shift in a high-precision survey in 2005, is the error
                noticeable?}
        \end{enumerate}
      \end{ideia}
    \end{column}
    \begin{column}{0.46\textwidth}
      \onslide<2->{
      \begin{regra}[title={\tr{Solução passo a passo}{Step-by-step Solution}}]
        \scriptsize\raggedright
        (a) $\Delta s = v \times \Delta t$\\[2pt]
        $\phantom{(a)} = \num{1.2}\,\mathrm{cm/ano} \times \num{10}\,\mathrm{anos} = \bnum{\num{12}\,cm}$.
        \medskip

        (b) \tr{\textbf{Sim, é enorme!} Para topografia comum de decímetro talvez passe, mas para
        controle geodésico e engenharia (\num{1} a \num{2}\,cm), \bnum{\num{12}\,cm} destrói a
        consistência do ajustamento!}{\textbf{Yes, it is huge!} For coarse decimeter mapping it
        might pass, but for geodetic control and engineering (\num{1} to \num{2}\,cm),
        \bnum{\num{12}\,cm} ruins the adjustment!}
      \end{regra}
      }
    \end{column}
  \end{columns}
\end{frame}

\begin{frame}{\tr{Para levar para casa (Aula 2)}{Takeaways (Class 2)}}
  \begin{enumerate}
    \setlength{\itemsep}{4pt}
    \scriptsize
    \item \tr{\textbf{Assunção 1993} foi o ponto de virada: abandonou-se a ilusão dos datums
          topocêntricos nacionais isolados em favor de um referencial geocêntrico continental.}{
          \textbf{Asunción 1993} was the turning point: isolated national topocentric datums were
          abandoned in favor of a continental geocentric reference frame.}
    \item \tr{O SIRGAS adotou o \textbf{ITRS} como sistema, o elipsoide \textbf{GRS80} e a tecnologia
          \textbf{GPS} em estrutura hierárquica por camadas.}{SIRGAS adopted \textbf{ITRS} as system,
          the \textbf{GRS80} ellipsoid, and \textbf{GPS} technology in a layered hierarchy.}
    \item \tr{\textbf{SIRGAS95} materializou o primeiro esqueleto com \bnum{\num{58}} estações
          (11 no Brasil) e precisão sub-centimétrica, amarradas ao ITRF94 na época \bnum{1995,4}.}{
          \textbf{SIRGAS95} materialized the first backbone with \bnum{\num{58}} stations (11 in Brazil)
          and sub-centimeter accuracy, tied to ITRF94 at epoch \bnum{1995.4}.}
    \item \tr{A grande lição foi que a \bnum{Terra é viva}: a crosta se deforma por tectônica de placas
          e sismos, exigindo modelos de velocidade e campanhas periódicas.}{The major lesson was that
          the \bnum{Earth is alive}: the crust deforms via plate tectonics and earthquakes, requiring
          velocity models and periodic campaigns.}
  \end{enumerate}
  \begin{ideia}
    \scriptsize \tr{\textbf{Na próxima aula:} a grandiosa campanha de maio de 2000, com 184 estações,
    e a matemática do ajustamento que gerou o SIRGAS2000 na época 2000,4.}{\textbf{Next class:} the
    massive May 2000 campaign with 184 stations, and the adjustment mathematics that generated
    SIRGAS2000 at epoch 2000.4.}
  \end{ideia}
\end{frame}

% ==================================================================
% AULA 3 — A REALIZAÇÃO: MAIO DE 2000 E O CÁLCULO
% ==================================================================
\section{\tr{Aula 3: A realização (maio de 2000)}{Class 3: The Realization (May 2000)}}

\begin{frame}{\tr{O gancho da Aula 3: da floresta ao cálculo}{Class 3 Hook: From the Rainforest to Computation}}
  \begin{ideia}
    \tr{\textbf{A pergunta desta aula:} como \num{184} receptores GPS espalhados por selvas, desertos,
    ilhas e capitais de dois continentes são transformados num \bnum{único conjunto de coordenadas
    milimétricas}?}{\textbf{The question for this class:} how do \num{184} GPS receivers scattered
    across rainforests, deserts, islands, and capitals in two continents get turned into a
    \bnum{single set of millimeter-accurate coordinates}?}
  \end{ideia}
  \vspace{0.2cm}
  \begin{itemize}
    \setlength{\itemsep}{5pt}
    \item \tr{Na Aula 2 vimos o projeto e a campanha piloto SIRGAS95 com \num{58} estações.}{In
          Class 2 we examined the design and the pilot SIRGAS95 campaign with \num{58} stations.}
    \item \tr{Em maio de 2000, o continente realizou uma mobilização científica sem precedentes:
          \bnum{\num{184} estações rastreando simultaneamente dia e noite}.}{In May 2000, the
          continent mounted an unprecedented scientific mobilization: \bnum{\num{184} stations
          tracking simultaneously day and night}.}
    \item \tr{Vamos entender a logística, a física das duplas diferenças e a matemática do
          \textbf{ajustamento por mínimos quadrados com injunções}.}{We will discover the logistics,
          the double-difference physics, and the mathematics of \textbf{least squares adjustment
          with constraints}.}
  \end{itemize}
\end{frame}

\begin{frame}{\tr{A Campanha de Maio de 2000: o salto de escala}{The May 2000 Campaign: The Leap in Scale}}
  \begin{columns}[onlytextwidth]
    \begin{column}{0.50\textwidth}
      \begin{tcolorbox}[colback=mTeal!8, colframe=mTeal, title={\tr{A Operação Continental}{The Continental Operation}}]
        \scriptsize
        \begin{itemize}
          \setlength{\itemsep}{3pt}
          \item \textbf{\tr{Datas:}{Dates:}} \bnum{10 a 19 de maio de 2000} (\num{10} dias consecutivos).
          \item \textbf{\tr{Sessões:}{Sessions:}} \num{24} horas contínuas de rastreio diário com taxa de \num{30} segundos.
          \item \textbf{\tr{Cobertura:}{Coverage:}} América do Sul, América Central e Caribe.
          \item \textbf{\tr{Total de estações:}{Total stations:}} \bnum{\num{184}} vértices geodésicos ocupados.
          \item \textbf{\tr{Estações no Brasil:}{Stations in Brazil:}} \bnum{\num{21}} estações continentais e insulares.
        \end{itemize}
      \end{tcolorbox}
    \end{column}
    \begin{column}{0.46\textwidth}
      \begin{itemize}
        \setlength{\itemsep}{4pt}
        \scriptsize
        \item \tr{\textbf{Três metas principais em 2000:}}{\textbf{Three primary goals in 2000:}}
              \begin{enumerate}
                \tiny
                \item \tr{\textbf{Reocupar estações de 1995:} medir o vetor de velocidade e deformação da crosta.}{\textbf{Reoccupy 1995 stations:} measure crustal velocity and deformation vectors.}
                \item \tr{\textbf{Conectar os marégrafos:} amarrar os pontos de origem altimétrica das redes de nivelamento ao elipsoide geocêntrico.}{\textbf{Tie tide gauges:} connect vertical datum origins to the geocentric ellipsoid.}
                \item \tr{\textbf{Densificar a rede:} fornecer vértices de controle a menos de \num{500}\,km de qualquer região desenvolvida.}{\textbf{Densify the network:} supply control marks within \num{500}\,km of any developed region.}
              \end{enumerate}
      \end{itemize}
    \end{column}
  \end{columns}
\end{frame}

\begin{frame}{\tr{A ``Foto de Formatura'': por que medir ao mesmo tempo?}{The ``Graduation Photo'': Why Measure Simultaneously?}}
  \begin{columns}[onlytextwidth]
    \begin{column}{0.54\textwidth}
      \begin{itemize}
        \setlength{\itemsep}{4pt}
        \scriptsize
        \item \tr{\textbf{A analogia escolar:} se cada formando tirar uma foto sozinho em dias e
              estúdios diferentes, como comparar a altura exata entre eles?}{\textbf{The school analogy:}
              if every graduate takes a photo alone on different days, how can you verify their exact relative height?}
        \item \tr{Na foto de formatura, \bnum{todos sobem no tablado juntos}: a mesma lente, a mesma
              iluminação e o mesmo instante de tempo!}{In a class photo, \bnum{everyone steps on the
              risers together}: the same lens, the same lighting, the exact same moment!}
        \item \tr{Na geodésia por satélite, os satélites GPS orbitam a $\approx \num{20200}$\,km de
              altitude a \bnum{\num{3.9}\,km/s}.}{In satellite geodesy, GPS satellites orbit at
              $\approx \num{20200}$\,km altitude traveling at \bnum{\num{3.9}\,km/s}.}
        \item \tr{Eles funcionam como uma \textbf{régua espacial em movimento rápido}: só servem de ponte
              se duas estações no chão olharem para eles \alert{no exato mesmo microsegundo}!}{\textbf{They
              act as a fast-moving space ruler}: they bridge two ground stations only when both track
              them at the \alert{exact same microsecond}!}
      \end{itemize}
    \end{column}
    \begin{column}{0.42\textwidth}
      \begin{regra}
        \scriptsize
        \tr{\textbf{O princípio do rastreio comum:}\\[3pt]
        A atmosfera e o relógio atômico do satélite variam continuamente.\\[3pt]
        Se a estação em Manaus e a estação em Curitiba registram a fase do mesmo satélite no mesmo
        instante, o erro do relógio do satélite \bnum{é idêntico} para as duas e desaparece na
        subtração!}{\textbf{Principle of common tracking:}\\[3pt]
        The atmosphere and the satellite's atomic clock drift continuously.\\[3pt]
        If Manaus and Curitiba record carrier phase from the same satellite at the exact same instant,
        the satellite clock error \bnum{is identical} for both and cancels out in subtraction!}
      \end{regra}
    \end{column}
  \end{columns}
\end{frame}

\begin{frame}{\tr{O milagre das duplas diferenças de fase}{The Double-Difference Phase Miracle}}
  \begin{columns}[onlytextwidth]
    \begin{column}{0.50\textwidth}
      \scriptsize
      \tr{A equação da fase observada ($\Phi$) entre a estação $i$ e o satélite $k$:}{The carrier phase
      equation ($\Phi$) between receiver $i$ and satellite $k$:}
      \begin{equation*}
        \Phi_i^k = \rho_i^k + c(\delta t_i - \delta t^k) - I_i^k + T_i^k + \lambda N_i^k + \varepsilon
      \end{equation*}
      \begin{enumerate}
        \setlength{\itemsep}{2pt}
        \tiny
        \item \textbf{\tr{Primeira diferença (entre estações):}{First difference (between receivers):}}
              $\Delta \Phi_{ij}^k = \Phi_j^k - \Phi_i^k$\\[1pt]
              $\implies$ \bnum{\tr{elimina o erro do relógio do satélite}{cancels satellite clock error}} $\delta t^k$!
        \item \textbf{\tr{Segunda diferença (entre satélites):}{Second difference (between satellites):}}
              $\nabla\Delta \Phi_{ij}^{kl} = \Delta \Phi_{ij}^l - \Delta \Phi_{ij}^k$\\[1pt]
              $\implies$ \bnum{\tr{elimina o erro do relógio do receptor}{cancels receiver clock error}} $\delta t_i, \delta t_j$!
      \end{enumerate}
    \end{column}
    \begin{column}{0.46\textwidth}
      \centering
      \begin{tikzpicture}[scale=0.85]
        \node (satA) at (-1.3,2.4) [caixa, fill=mAmber!20] {\tiny Satélite $k$};
        \node (satB) at (1.3,2.4) [caixa, fill=mAmber!20] {\tiny Satélite $l$};
        \pic at (-1.5,0) {antena};
        \pic at (1.5,0) {antena};
        \draw[thick] (-2,0) -- (2,0);
        \node at (-1.5,-0.3) [font=\tiny] {\tr{Estação $i$}{Station $i$}};
        \node at (1.5,-0.3) [font=\tiny] {\tr{Estação $j$}{Station $j$}};
        \draw[seta, dashed, mTeal] (satA) -- (-1.5,0.22);
        \draw[seta, dashed, mTeal] (satA) -- (1.5,0.22);
        \draw[seta, dashed, mIndigo] (satB) -- (-1.5,0.22);
        \draw[seta, dashed, mIndigo] (satB) -- (1.5,0.22);
        \draw[seta2, thick, mRose] (-1.3,0) -- (1.3,0) node[midway, below, font=\tiny] {$\vec{X}_j - \vec{X}_i$};
      \end{tikzpicture}\\[3pt]
      \tiny \tr{Restam apenas a geometria pura do vetor e a ambiguidade inteira $N$.}{Only pure
      geometry of the baseline vector and integer ambiguity $N$ remain.}
    \end{column}
  \end{columns}
\end{frame}

\begin{frame}{\tr{O que cancela e o que NÃO cancela?}{What Cancels and What Does NOT?}}
  \begin{columns}[onlytextwidth]
    \begin{column}{0.48\textwidth}
      \begin{tcolorbox}[colback=mTeal!8, colframe=mTeal, title={\tr{O que se cancela / modela}{What cancels / is modeled}}]
        \scriptsize
        \begin{itemize}
          \setlength{\itemsep}{2pt}
          \item \textbf{\tr{Relógios:}{Clocks:}} erro de relógio do satélite e do receptor (eliminados pelas diferenças).
          \item \textbf{\tr{Órbitas:}{Orbits:}} produtos orbitais precisos finais do IGS (\num{1} a \num{2}\,cm no espaço).
          \item \textbf{\tr{Ionosfera:}{Ionosphere:}} combinação livre da ionosfera L3 / L4 (combinação linear das frequências L1 e L2).
        \end{itemize}
      \end{tcolorbox}
    \end{column}
    \begin{column}{0.48\textwidth}
      \begin{tcolorbox}[colback=mRose!8, colframe=mRose, title={\tr{O que não se cancela}{What does not cancel}}]
        \scriptsize
        \begin{itemize}
          \setlength{\itemsep}{2pt}
          \item \textbf{\tr{Troposfera local:}{Local troposphere:}} vapor d'água na baixa atmosfera (exige estimar parâmetros de zênite a cada hora).
          \item \textbf{\tr{Multicaminho:}{Multipath:}} reflexão do sinal em construções, solo ou veículos vizinhos.
          \item \textbf{\tr{Centro de fase da antena:}{Antenna phase center:}} variações mecânicas dependentes do ângulo zenital.
        \end{itemize}
      \end{tcolorbox}
    \end{column}
  \end{columns}
  \vspace{0.25cm}
  \begin{cuidado}
    \scriptsize \tr{\textbf{Atenção:} Para linhas de base de \num{2000}\,km, a atmosfera sobre Manaus
    é completamente diferente da atmosfera sobre Curitiba. O atraso troposférico precisa ser
    estimado estação por estação!}{\textbf{Notice:} Over \num{2000}\,km baselines, the atmosphere
    over Manaus is completely different from Curitiba. Tropospheric delay must be solved per station!}
  \end{cuidado}
\end{frame}

\begin{frame}{\tr{Três Centros de Análise: ciência sem viés}{Three Analysis Centers: Unbiased Science}}
  \begin{columns}[onlytextwidth]
    \begin{column}{0.54\textwidth}
      \begin{itemize}
        \setlength{\itemsep}{4pt}
        \scriptsize
        \item \tr{Assim como em 1995, os dados brutos de maio de 2000 foram processados por
              \bnum{três centros de processamento independentes}:}{As in 1995, raw tracking data
              from May 2000 were processed by \bnum{three independent analysis centers}:}
              \begin{enumerate}
                \tiny
                \item \textbf{DGFI} (\emph{Deutsches Geodätisches Forschungsinstitut}, Munique) ---
                      Bernese GPS Software 4.2;
                \item \textbf{IBGE} (Rio de Janeiro, Brasil) --- Bernese GPS Software 4.0;
                \item \textbf{BEK} (\emph{Bayerische Kommission für die Internationale Erdmessung}) ---
                      GIPSY/OASIS II.
              \end{enumerate}
        \item \tr{Repare: \textbf{duas famílias matemáticas} inteiramente distintas (Bernese por duplas
              diferenças; GIPSY por posicionamento por ponto preciso com filtro de Kalman).}{Notice:
              \textbf{two entirely different mathematical methodologies} (Bernese via double differences;
              GIPSY via PPP with Kalman filters).}
      \end{itemize}
    \end{column}
    \begin{column}{0.42\textwidth}
      \begin{ideia}
        \scriptsize
        \tr{\textbf{Por que múltiplos centros?}\\[3pt]
        Se um único centro calcula, um erro sutil no algoritmo ou na modelagem da maré terrestre pode
        contaminar toda a cartografia de um continente.\\[3pt]
        A comparação cega entre soluções independentes garante que o resultado é \bnum{física pura e
        reprodutível}, e não artefato computacional.}{\textbf{Why multiple centers?}\\[3pt]
        If a single center computes, a subtle code bug or solid Earth tide model error could taint
        a continent's mapping.\\[3pt]
        Blind comparison across independent solutions guarantees the output is \bnum{pure,
        reproducible physics}, not an algorithmic artifact.}
      \end{ideia}
    \end{column}
  \end{columns}
\end{frame}

\begin{frame}{\tr{O formato SINEX: combinando equações sem deformação}{The SINEX Format: Combining Equations Without Distortion}}
  \begin{columns}[onlytextwidth]
    \begin{column}{0.52\textwidth}
      \begin{itemize}
        \setlength{\itemsep}{4pt}
        \scriptsize
        \item \tr{Cada centro de análise não publicou coordenadas finais prontas, mas sim arquivos
              no formato neutro \textbf{SINEX} (\emph{Solution INdependent EXchange}).}{Each center
              published not final rigid coordinates, but neutral \textbf{SINEX} files (\emph{Solution
              INdependent EXchange}).}
        \item \tr{O arquivo SINEX contém:}{The SINEX file contains:}
              \begin{itemize}
                \tiny
                \item A matriz de equações normais reduzida ($N = A^T P A$);
                \item O vetor de termos independentes ($u = A^T P l$);
                \item A matriz de variância-covariância completa entre todas as estações.
              \end{itemize}
        \item \tr{Isso permitiu combinar as soluções como se fossem um único mega-ajustamento, sem
              arredondamentos nem perda de correlação entre os marcos.}{This enabled combining solutions
              as a unified mega-adjustment without rounding or covariance loss.}
      \end{itemize}
    \end{column}
    \begin{column}{0.44\textwidth}
      \begin{regra}
        \scriptsize
        \tr{\textbf{A combinação por Mínimos Quadrados:}\\[3pt]
        Dadas as normais $N_1, N_2, N_3$ dos três centros com fatores de escala de variância
        $\sigma_1^2, \sigma_2^2, \sigma_3^2$:\\[4pt]
        \centering
        $N_{\text{comb}} = \sum_{k=1}^3 \dfrac{1}{\sigma_k^2} N_k$\\[4pt]
        $u_{\text{comb}} = \sum_{k=1}^3 \dfrac{1}{\sigma_k^2} u_k$\\[4pt]
        \raggedright
        A combinação é rigorosamente ótima no sentido dos mínimos quadrados!}{
        \textbf{Least Squares Combination:}\\[3pt]
        Given normal systems $N_1, N_2, N_3$ from the three centers with variance factors
        $\sigma_1^2, \sigma_2^2, \sigma_3^2$:\\[4pt]
        \centering
        $N_{\text{comb}} = \sum_{k=1}^3 \dfrac{1}{\sigma_k^2} N_k$\\[4pt]
        $u_{\text{comb}} = \sum_{k=1}^3 \dfrac{1}{\sigma_k^2} u_k$\\[4pt]
        \raggedright
        The combination is strictly optimal in the least-squares sense!}
      \end{regra}
    \end{column}
  \end{columns}
\end{frame}

\begin{frame}{\tr{A Rede Livre e o Defeito de Posto}{The Free Network and Rank Defect}}
  \begin{columns}[onlytextwidth]
    \begin{column}{0.52\textwidth}
      \begin{itemize}
        \setlength{\itemsep}{4pt}
        \scriptsize
        \item \tr{A rede GPS combinada determina os comprimentos relativos e os ângulos entre as
              \num{184} estações com precisão de \bnum{\num{2} a \num{4}\,mm}.}{The combined GPS
              network determines relative baseline lengths and angles among \num{184} stations with
              \bnum{\num{2} to \num{4}\,mm} precision.}
        \item \tr{No entanto, essa geometria perfeita está \textbf{solta no espaço}:}{However, this
              perfect geometry is \textbf{unattached in space}:}
              \begin{itemize}
                \tiny
                \item Ela não sabe onde fica a origem $(0, 0, 0)$ do geocentro (defeito de translação de rank = 3);
                \item A matriz de equações normais $N$ é \alert{singular} ($\det N = 0$) e não tem inversa direta!
              \end{itemize}
        \item \tr{Para transformar essa forma geométrica solta num referencial oficial, precisamos
              \bnum{ancorar a rede no ITRS}.}{To convert this floating geometry into an official frame,
              we must \bnum{anchor the network to ITRS}.}
      \end{itemize}
    \end{column}
    \begin{column}{0.44\textwidth}
      \centering
      \begin{tikzpicture}[scale=0.85]
        \draw[thick, corFrame, fill=corFrame!10] (-1.5,-0.5) -- (1.5,-0.8) -- (1.8,1.2) -- (-0.8,1.5) -- cycle;
        \fill[corFrame] (-1.5,-0.5) circle (2pt) node[below, font=\tiny] {$A$};
        \fill[corFrame] (1.5,-0.8) circle (2pt) node[below, font=\tiny] {$B$};
        \fill[corFrame] (1.8,1.2) circle (2pt) node[above, font=\tiny] {$C$};
        \fill[corFrame] (-0.8,1.5) circle (2pt) node[above, font=\tiny] {$D$};
        \node at (0.2,0.3) [font=\tiny, text=corFrame] {\tr{Forma interna rígida}{Rigid internal shape}};
        \draw[seta2, corErro, thick] (0.2,0.6) -- (1.6,2.0) node[midway, above right, font=\tiny] {\tr{Flutua no espaço $(\Delta X, \Delta Y, \Delta Z)$}{Floats in space}};
      \end{tikzpicture}\\[2pt]
      \tiny \tr{Sem âncoras, a rede pode transladar em bloco sem alterar nenhuma observação.}{Without
      anchors, the entire network can shift without altering any internal observation.}
    \end{column}
  \end{columns}
\end{frame}

\begin{frame}{\tr{Ajustamento com Injunções: o modelo matemático}{Adjustment with Constraints: The Mathematical Model}}
  \begin{regra}
    \scriptsize
    \tr{\textbf{Como pregar a rede no referencial global (ITRF2000):}\\[2pt]
    Adicionam-se equações de observação para as coordenadas conhecidas de estações de referência
    (estações do IGS):\\[3pt]
    \centering
    $C \hat{x} = w + v_C, \quad \text{com matriz de pesos } P_C = \Sigma_C^{-1}$\\[3pt]
    \raggedright
    O sistema de equações normais com injunções ponderadas torna-se:\\[3pt]
    \centering
    $\bnum{(N + C^T P_C C) \hat{x} = u + C^T P_C w}$}{\textbf{How to anchor the network to the
    global frame (ITRF2000):}\\[2pt]
    We append observation equations for known coordinates of reference stations (IGS stations):\\[3pt]
    \centering
    $C \hat{x} = w + v_C, \quad \text{with weight matrix } P_C = \Sigma_C^{-1}$\\[3pt]
    \raggedright
    The normal equation system with weighted constraints becomes:\\[3pt]
    \centering
    $\bnum{(N + C^T P_C C) \hat{x} = u + C^T P_C w}$}
  \end{regra}
  \vspace{0.15cm}
  \begin{columns}[onlytextwidth]
    \begin{column}{0.48\textwidth}
      \begin{tcolorbox}[colback=mStoneLight, colframe=mStone, title={\tr{Injunção Mínima}{Minimal Constraint}}]
        \tiny \tr{Fixa apenas 3 parâmetros de translação (condição de Helmert). A rede interna não
        sofre nenhuma tensão artificial.}{Fixes only 3 translation parameters (Helmert condition).
        The internal network experiences zero artificial stress.}
      \end{tcolorbox}
    \end{column}
    \begin{column}{0.48\textwidth}
      \begin{tcolorbox}[colback=mTeal!8, colframe=mTeal, title={\tr{Injunção Ponderada no ITRF}{Weighted Constraint to ITRF}}]
        \tiny \tr{As coordenadas IGS entram com suas incertezas oficiais ($\Sigma_C$). A rede é
        amarrada ao geocentro com máxima fidelidade física.}{IGS coordinates enter with official
        uncertainties ($\Sigma_C$). The network ties to the geocenter with optimal fidelity.}
      \end{tcolorbox}
    \end{column}
  \end{columns}
\end{frame}

\begin{frame}{\tr{Por que a época é 2000,4?}{Why is the Epoch 2000.4?}}
  \begin{columns}[onlytextwidth]
    \begin{column}{0.54\textwidth}
      \begin{itemize}
        \setlength{\itemsep}{4pt}
        \scriptsize
        \item \tr{A campanha ocorreu entre os dias \textbf{10 e 19 de maio de 2000}.}{The campaign
              took place between \textbf{May 10 and 19, 2000}.}
        \item \tr{No calendário geodésico, o dia médio da campanha é o dia do ano \bnum{\num{145}}
              (\num{24} de maio):}{In the geodetic calendar, the campaign mean day is day of year
              \bnum{\num{145}} (May 24):}
              \begin{equation*}
                t = 2000 + \frac{145}{365{,}25} = \bnum{\num{2000.4}}
              \end{equation*}
        \item \tr{O sistema de referência adotado foi o \textbf{ITRF2000}, propagado para a época
              \bnum{2000,4}.}{The global reference frame adopted was \textbf{ITRF2000}, propagated
              to epoch \bnum{2000.4}.}
        \item \tr{As coordenadas resultantes são exatamente as posições das estações \textbf{naquele dia
              de maio de 2000}!}{The resulting coordinates represent the exact positions of stations
              \textbf{on that day of May 2000}!}
      \end{itemize}
    \end{column}
    \begin{column}{0.42\textwidth}
      \begin{cuidado}
        \scriptsize
        \tr{\textbf{O conceito de Frame Congelado:}\\[3pt]
        O SIRGAS2000 nasce como um \emph{frame estático}.\\[3pt]
        Ao contrário do ITRF contínuo (cujas coordenadas mudam todos os dias pela tectônica de
        placas), o SIRGAS2000 foi \bnum{congelado na época 2000,4}.\\[3pt]
        Isso era indispensável para que as cartas e cadastros de propriedades rurais e urbanas não
        precisassem ser redesenhados a cada semana!}{\textbf{The Frozen Frame Concept:}\\[3pt]
        SIRGAS2000 was born as a \emph{static frame}.\\[3pt]
        Unlike continuous ITRF (whose coordinates drift daily due to plate tectonics), SIRGAS2000 was
        \bnum{frozen at epoch 2000.4}.\\[3pt]
        This was essential so cadastral boundaries and national map sheets would not have to be
        redrawn every week!}
      \end{cuidado}
    \end{column}
  \end{columns}
\end{frame}

\begin{frame}{\tr{As 21 estações brasileiras na Campanha de 2000}{The 21 Brazilian Stations in the 2000 Campaign}}
  \begin{columns}[onlytextwidth]
    \begin{column}{0.54\textwidth}
      \scriptsize
      \tr{Os \num{21} vértices fundamentais que ancoraram o novo SGB no território:}{The \num{21}
      fundamental stations anchoring the new SGB across Brazil:}
      \begin{itemize}
        \setlength{\itemsep}{1.5pt}
        \tiny
        \item \textbf{Norte:} Belém (PA), Macapá (AP), Manaus (AM), Oiapoque (AP), Porto Velho (RO), Rio Branco (AC).
        \item \textbf{Nordeste:} Fortaleza (CE), Imperatriz (MA), Natal (RN), Petrolina (PE), Salvador (BA), São Luís (MA).
        \item \textbf{Centro-Oeste:} Brasília (DF), Campo Grande (MS), Cuiabá (MT).
        \item \textbf{Sudeste:} Cananéia (SP), Presidente Prudente (SP), Rio de Janeiro (RJ), Viçosa (MG).
        \item \textbf{Sul:} Curitiba (PR), Santa Maria (RS).
      \end{itemize}
      \vspace{0.1cm}
      \tr{No ajuste da rede planimétrica de 2006, \bnum{20 dessas estações continentais mais SMAR
      (Santa Maria)} serviram como as \num{21} equações de injunção que fixaram todo o país! [B]}{
      In the 2006 network adjustment, \bnum{20 continental sites plus SMAR (Santa Maria)} served
      as the \num{21} constraint equations locking the entire country! [B]}
    \end{column}
    \begin{column}{0.42\textwidth}
      \begin{ideia}
        \scriptsize
        \tr{\textbf{A curiosidade de Cananéia:}\\[3pt]
        A estação CANA (Cananéia/SP) foi uma das 21 ocupadas na campanha de maio de 2000.\\[3pt]
        No entanto, o marco foi destruído posteriormente por obras de engenharia local.\\[3pt]
        Por isso, no cálculo de 2006, o IBGE substituiu Cananéia por \textbf{SMAR} (Santa Maria),
        mantendo exatamente \bnum{21 estações de injunção} na rede nacional! [B]}{\textbf{The Cananéia
        curiosity:}\\[3pt]
        Station CANA was one of 21 occupied in May 2000.\\[3pt]
        However, the mark was subsequently destroyed by local construction.\\[3pt]
        Therefore, in the 2006 calculation, IBGE replaced Cananéia with \textbf{SMAR} (Santa Maria),
        retaining exactly \bnum{21 constraint stations} across the nation! [B]}
      \end{ideia}
    \end{column}
  \end{columns}
\end{frame}

\begin{frame}{\tr{A integração dos marégrafos: a promessa cumprida do GT III}{Tide Gauge Integration: Fulfilling the WG III Promise}}
  \begin{columns}[onlytextwidth]
    \begin{column}{0.50\textwidth}
      \begin{itemize}
        \setlength{\itemsep}{4pt}
        \scriptsize
        \item \tr{Lembre-se: em 1997, na Assembleia da IAG no Rio, o SIRGAS criou o \textbf{GT III
              (Datum Vertical)}.}{Recall: in 1997 at the IAG Assembly in Rio, SIRGAS formed
              \textbf{WG III (Vertical Datum)}.}
        \item \tr{Em maio de 2000, pela primeira vez na América Latina, antenas GPS ocuparam marcos
              vizinhos aos principais marégrafos do continente:}{In May 2000, for the first time in
              Latin America, GPS antennas occupied marks adjacent to primary continental tide gauges:}
              \begin{itemize}
                \tiny
                \item Marégrafo de \textbf{Imbituba (SC)} --- a origem altimétrica do SGB;
                \item Marégrafo de \textbf{Santana (AP)} --- a origem do datum vertical amapaense;
                \item Marégrafos na Argentina, Chile, Colômbia, Venezuela e Uruguai.
              \end{itemize}
      \end{itemize}
    \end{column}
    \begin{column}{0.46\textwidth}
      \begin{regra}
        \scriptsize
        \tr{\textbf{O significado revolucionário:}\\[3pt]
        A altimetria clássica dá altitude sobre o mar ($H$). O GPS dá altura sobre o elipsoide ($h$).\\[3pt]
        Ao medir o marégrafo com GPS, descobriu-se com precisão milimétrica a \bnum{ondulação geoidal
        $N = h - H$} no ponto de partida das redes de nivelamento!\\[3pt]
        Essa ponte foi o primeiro passo para o modelo de quase-geoide e a unificação altimétrica (tema
        da nossa Parte 2!).}{\textbf{Revolutionary significance:}\\[3pt]
        Leveling yields height above sea level ($H$). GPS yields ellipsoidal height ($h$).\\[3pt]
        Surveying the tide gauge with GPS revealed with millimeter accuracy the \bnum{geoidal
        undulation $N = h - H$} at the leveling datum origin!\\[3pt]
        This bridge enabled modern geoid modeling and vertical unification (the topic of our Part 2!).}
      \end{regra}
    \end{column}
  \end{columns}
\end{frame}

\begin{frame}{\tr{Pare e pense: por que o Brasil não adotou o SIRGAS logo em 2000?}{Stop and Think: Why Did Brazil Wait After 2000?}}
  \begin{armadilha}
    \tr{\textbf{A dúvida comum:} ``Se a campanha continental terminou em maio de 2000 e o cálculo
    ficou pronto pouco depois, por que o Brasil só oficializou o SIRGAS2000 em 2005 e só publicou o
    reajuste da rede em 2006?''}{\textbf{Common doubt:} ``If the continental campaign finished in
    May 2000, why did Brazil wait until 2005 to decree SIRGAS2000 and 2006 to publish the national
    network adjustment?''}
  \end{armadilha}
  \vspace{0.15cm}
  \begin{columns}[onlytextwidth]
    \begin{column}{0.48\textwidth}
      \begin{tcolorbox}[colback=mStoneLight, colframe=mStone, title={\tr{O abismo entre o esqueleto e a rede}{Gap between backbone and network}}]
        \scriptsize
        \begin{itemize}
          \setlength{\itemsep}{2pt}
          \item \tr{Ter \num{21} estações milimétricas é apenas o ``esqueleto''.}{Having \num{21} millimeter stations is merely the ``backbone''.}
          \item \tr{O Brasil real tinha mais de \bnum{\num{4700} marcos clássicos} de concreto em morros e fazendas amarrados ao SAD69.}{Brazil had over \bnum{\num{4700} classical marks} on hilltops tied to SAD69.}
        \end{itemize}
      \end{tcolorbox}
    \end{column}
    \begin{column}{0.48\textwidth}
      \begin{tcolorbox}[colback=mTeal!8, colframe=mTeal, title={\tr{O trabalho hercúleo (2000--2005)}{The Herculean task (2000--2005)}}]
        \scriptsize
        \begin{itemize}
          \setlength{\itemsep}{2pt}
          \item \tr{Criou-se o PMRG (2000) e realizaram-se \bnum{\num{12} campanhas GPS adicionais} em 2004--2005 ligando a rede clássica.}{PMRG was created (2000) and \bnum{\num{12} GPS campaigns} ran in 2004--2005 tying classical marks.}
          \item \tr{Era preciso reajustar simultaneamente \num{6265} marcos sem quebrar as fronteiras estaduais! [B]}{All \num{6265} marks had to be re-adjusted simultaneously without border seams! [B]}
        \end{itemize}
      \end{tcolorbox}
    \end{column}
  \end{columns}
\end{frame}

\begin{frame}{\tr{Sua vez: a lógica do peso da injunção}{Your Turn: The Logic of Constraint Weights}}
  \begin{columns}[T]
    \begin{column}{0.50\textwidth}
      \begin{ideia}[title={\tr{Problema}{Problem}}]
        \scriptsize
        \tr{No ajuste de 2006 da rede planimétrica, o IBGE introduziu as \num{21} estações SIRGAS com
        injunção ponderada, usando desvio-padrão dividido por \num{10}: $\sigma_{\text{inj}} = \sigma / 10$.}{
        In the 2006 network adjustment, IBGE constrained the \num{21} SIRGAS stations
        using a standard deviation divided by \num{10}: $\sigma_{\text{inj}} = \sigma / 10$.}
        \begin{enumerate}
          \setlength{\itemsep}{2pt}
          \item[(a)] \tr{Como o peso é proporcional a $1/\sigma^2$, por quanto o peso da estação foi
                multiplicado?}{Since weight is proportional to $1/\sigma^2$, by what factor was the station weight
                multiplied?}
          \item[(b)] \tr{Qual a consequência prática disso sobre as estações da rede clássica?}{What was the
                practical consequence on the classical network stations?}
        \end{enumerate}
      \end{ideia}
    \end{column}
    \begin{column}{0.46\textwidth}
      \begin{regra}[title={\tr{Respostas}{Answers}}]
        \scriptsize\raggedright
        (a) $P = \dfrac{1}{(\sigma/10)^2} = \bnum{\num{100}} \times P_0$.\\[2pt]
        \tr{O peso aumentou \bnum{\num{100} vezes}! [B]}{Weight increased by \bnum{\num{100} times}! [B]}
        \medskip

        (b) \tr{\textbf{Rigidez quase absoluta:} as \num{21} estações SIRGAS comportam-se como
        colunas de concreto inabaláveis. A rede clássica é forçada a se ajustar a elas, sem que os
        erros antigos desloquem o referencial geocêntrico!}{\textbf{Near-absolute stiffness:} the
        \num{21} SIRGAS stations act as unyielding pillars. The classical network is forced to fit
        them, preventing old distortions from pulling the geocentric frame!}
      \end{regra}
    \end{column}
  \end{columns}
\end{frame}

\begin{frame}{\tr{Para levar para casa (Aula 3)}{Takeaways (Class 3)}}
  \begin{enumerate}
    \setlength{\itemsep}{4pt}
    \scriptsize
    \item \tr{A campanha de \textbf{maio de 2000} ocupou simultaneamente \bnum{\num{184} estações}
          em dois continentes por \num{10} dias de \num{24} horas, incluindo \bnum{\num{21}} no Brasil.}{
          The \textbf{May 2000} campaign simultaneously occupied \bnum{\num{184} stations} across
          two continents for \num{10} days of \num{24} hours, including \bnum{\num{21}} in Brazil.}
    \item \tr{O rastreio simultâneo e as \textbf{duplas diferenças de fase} eliminaram os erros de
          relógio dos satélites e receptores, fornecendo precisão de milímetros.}{Simultaneous tracking
          and \textbf{carrier double differences} eliminated satellite and receiver clock errors,
          yielding millimeter precision.}
    \item \tr{Três centros independentes (DGFI, IBGE e BEK) processaram os dados; as equações normais
          em formato \textbf{SINEX} foram combinadas e ajustadas com \bnum{injunções ponderadas no ITRF2000}.}{
          Three independent centers (DGFI, IBGE, BEK) processed data; normal equations in \textbf{SINEX}
          format were combined and adjusted with \bnum{weighted constraints to ITRF2000}.}
    \item \tr{O frame nasceu congelado na época \bnum{2000,4}, amarrando também pela primeira vez os
          marégrafos das redes de nivelamento físico (Imbituba, Santana).}{The frame was frozen at epoch
          \bnum{2000.4}, also tying for the first time the physical leveling tide gauges (Imbituba, Santana).}
  \end{enumerate}
  \begin{ideia}
    \scriptsize \tr{\textbf{Na próxima aula:} a adoção legal no Brasil (Decreto 5.334/2005 e R.PR 1/2005),
    o reajuste da rede em 2006 com 6265 estações, os parâmetros de conversão e a vida útil do referencial.}{
    \textbf{Next class:} legal adoption in Brazil (Decree 5,334/2005 and R.PR 1/2005), the 2006 network
    adjustment with 6,265 stations, conversion parameters, and frame lifetime.}
  \end{ideia}
\end{frame}


% ==================================================================
% AULA 4 — ADOÇÃO, DENSIFICAÇÃO E VIDA ÚTIL
% ==================================================================
\section{\tr{Aula 4: Adoção, densificação e vida útil}{Class 4: Adoption, Densification and Lifetime}}

\begin{frame}{\tr{O gancho da Aula 4: o referencial congelado e o chão que anda}{Class 4 Hook: The Frozen Frame and the Moving Ground}}
  \begin{ideia}
    \tr{\textbf{A pergunta desta aula:} uma vez calculadas as \num{21} estações fundamentais em maio
    de 2000, como isso chega ao mapa do município, ao título de terra e ao receptor no campo? E se o
    referencial está congelado na época \bnum{2000,4}, por que a coordenada do receptor hoje não bate
    mais com a monografia oficial?}{\textbf{The question for this class:} once the \num{21}
    fundamental stations were computed in May 2000, how does that reach municipal maps, land titles,
    and field receivers? And if the frame was frozen at epoch \bnum{2000.4}, why does a modern
    receiver coordinate drift from the official monument datasheet?}
  \end{ideia}
  \vspace{0.2cm}
  \begin{itemize}
    \setlength{\itemsep}{5pt}
    \item \tr{Na Aula 3 vimos a campanha de maio de 2000 e o cálculo do esqueleto continental.}{In
          Class 3 we saw the May 2000 campaign and the computation of the continental backbone.}
    \item \tr{Nesta aula final da Parte 1, veremos como o IBGE transferiu essa exatidão para todo o
          território nacional, o marco legal da adoção e a convivência com o movimento das placas.}{In
          this concluding class of Part 1, we will see how IBGE densified this accuracy across Brazil,
          the legal adoption framework, and living with tectonic plate motion.}
  \end{itemize}
\end{frame}

\begin{frame}{\tr{O desafio da densificação: um país continental}{The Densification Challenge: A Continental Nation}}
  \begin{columns}[onlytextwidth]
    \begin{column}{0.52\textwidth}
      \begin{itemize}
        \setlength{\itemsep}{4pt}
        \scriptsize
        \item \tr{O Brasil tem mais de \bnum{\num{8.5} milhões de $\mathrm{km}^2$}. Vinte e uma
              estações espaçadas por milhares de quilômetros não bastam para o dia a dia da agrimensura.}{
              Brazil covers over \bnum{\num{8.5} million $\mathrm{km}^2$}. Twenty-one stations
              separated by thousands of kilometers cannot directly support daily land surveying.}
        \item \tr{Em 2000, o IBGE formalizou o \textbf{PMRG} (\emph{Programa de Modernização da Rede
              Geodésica}), com duas frentes de ação:}{In 2000, IBGE formalized the \textbf{PMRG}
              (\emph{Geodetic Network Modernization Program}), with two action fronts:}
              \begin{enumerate}
                \tiny
                \item \tr{\textbf{Rede Ativa (RBMC):} estações GNSS contínuas conectadas à internet transmitindo dados 24h por dia.}{\textbf{Active Network (RBMC):} continuous GNSS stations connected to the internet streaming data 24/7.}
                \item \tr{\textbf{Reajuste da Rede Passiva:} conectar os milhares de pilares clássicos de concreto ao novo referencial SIRGAS2000.}{\textbf{Passive Network Re-adjustment:} tie the thousands of classical concrete monuments to SIRGAS2000.}
              \end{enumerate}
      \end{itemize}
    \end{column}
    \begin{column}{0.44\textwidth}
      \begin{tcolorbox}[colback=mIndigo!8, colframe=mIndigo, title={\tr{O nascimento do PMRG}{Birth of PMRG}}]
        \scriptsize \tr{Instituído no 1.º Seminário sobre o Sistema Geodésico Brasileiro (Curitiba,
        2000), atendendo à recomendação do Congresso Brasileiro de Cartografia de 1999. [B]}{Established
        at the 1st Seminar on the Brazilian Geodetic System (Curitiba, 2000), fulfilling the
        recommendation of the 1999 Brazilian Cartography Congress. [B]}
      \end{tcolorbox}
    \end{column}
  \end{columns}
\end{frame}

\begin{frame}{\tr{A Rede Ativa: a RBMC}{The Active Network: RBMC}}
  \begin{columns}[onlytextwidth]
    \begin{column}{0.54\textwidth}
      \begin{itemize}
        \setlength{\itemsep}{4pt}
        \scriptsize
        \item \tr{A \textbf{RBMC} (\emph{Rede Brasileira de Monitoramento Contínuo dos Sistemas GNSS})
              é a espinha dorsal ativa do país.}{The \textbf{RBMC} (\emph{Brazilian Continuous
              Monitoring Network of GNSS Systems}) is the active backbone of Brazil.}
        \item \tr{\textbf{Histórico pioneiro:} inaugurada em \bnum{dezembro de 1996} com duas
              estações pioneiras:}{
              \textbf{Pioneering history:} inaugurated in \bnum{December 1996} with two pioneering
              stations:}
              \begin{itemize}
                \tiny
                \item Curitiba (UFPR, Paraná);
                \item Presidente Prudente (UNESP, São Paulo).
              \end{itemize}
        \item \tr{Hoje conta com mais de uma centena de estações com dados horários e diários em
              formato RINEX disponibilizados gratuitamente pelo IBGE.}{Today it comprises over one
              hundred stations providing hourly and daily RINEX files free of charge via IBGE.}
      \end{itemize}
    \end{column}
    \begin{column}{0.42\textwidth}
      \begin{ideia}
        \scriptsize
        \tr{\textbf{A revolução da rede ativa:}\\[3pt]
        Em vez de procurar um marco de concreto escondido na vegetação, o usuário rastreia com seu
        receptor e calcula o vetor de base relativo à estação RBMC mais próxima, cuja coordenada
        oficial é conhecida no milímetro!}{\textbf{The active network revolution:}\\[3pt]
        Instead of hunting for a concrete monument hidden in brush, surveyors track with their own
        receiver and compute a baseline relative to the nearest RBMC station, whose official
        coordinate is known to the millimeter!}
      \end{ideia}
    \end{column}
  \end{columns}
\end{frame}

\begin{frame}{\tr{O Serviço IBGE-PPP: democratizando a alta precisão}{The IBGE-PPP Service: Democratizing High Precision}}
  \begin{columns}[onlytextwidth]
    \begin{column}{0.50\textwidth}
      \begin{itemize}
        \setlength{\itemsep}{4pt}
        \scriptsize
        \item \tr{O \textbf{IBGE-PPP} é um serviço online público e gratuito para processamento de
              dados GNSS pelo método de Posicionamento por Ponto Preciso.}{The \textbf{IBGE-PPP} is
              a free public online service for processing GNSS data using Precise Point Positioning.}
        \item \tr{Desenvolvido pelo IBGE em parceria com o \emph{Natural Resources Canada} (NRCan).}{
              Developed by IBGE in partnership with Natural Resources Canada (NRCan).}
        \item \tr{\textbf{Como funciona:}}{ \textbf{How it works:}}
              \begin{enumerate}
                \tiny
                \item O usuário envia um arquivo RINEX de rastreio estático pelo navegador;
                \item O servidor aplica órbitas e relógios precisos do IGS;
                \item O relatório PDF sai em minutos com coordenadas oficiais no \bnum{SIRGAS2000 na
                      época 2000,4}.
              \end{enumerate}
      \end{itemize}
    \end{column}
    \begin{column}{0.46\textwidth}
      \begin{regra}
        \scriptsize
        \tr{\textbf{Impacto no país:}\\[3pt]
        Qualquer cidadão, prefeitura ou pequeno proprietário com um receptor geodésico tem acesso
        direto à mesma qualidade geodésica do IBGE, sem softwares comerciais de dezenas de milhares
        de reais.}{\textbf{Impact across Brazil:}\\[3pt]
        Any citizen, municipality, or surveyor with a geodetic receiver gains direct access to
        IBGE's geodetic accuracy, without needing expensive commercial software packages.}
      \end{regra}
    \end{column}
  \end{columns}
\end{frame}

\begin{frame}{\tr{O Ajustamento da Rede Planimétrica de 2006}{The 2006 Planimetric Network Adjustment}}
  \begin{columns}[onlytextwidth]
    \begin{column}{0.52\textwidth}
      \begin{itemize}
        \setlength{\itemsep}{4pt}
        \scriptsize
        \item \tr{Em junho de \bnum{2006}, o IBGE concluiu o maior ajustamento por mínimos quadrados
              da história da América do Sul. [B]}{In June \bnum{2006}, IBGE concluded the largest
              least-squares adjustment in South American history. [B]}
        \item \tr{\textbf{Dimensão do ajustamento:}}{\textbf{Adjustment scale:}}
              \begin{itemize}
                \tiny
                \item \bnum{\num{6265}} estações ajustadas simultaneamente;
                \item \num{4774} marcos da rede clássica de triangulação/poligonação;
                \item \num{1600} estações ocupadas por GPS moderno.
              \end{itemize}
        \item \tr{\textbf{Total de observações (22.232 equações):}}{\textbf{Total observations (22,232 equations):}}
              \begin{itemize}
                \tiny
                \item \num{16913} direções horizontais;
                \item \num{1534} distâncias medidas;
                \item \num{389} azimutes de Laplace;
                \item \num{179} posições Doppler;
                \item \num{3196} diferenças de coordenadas GPS;
                \item \bnum{\num{21}} equações de injunção (estações SIRGAS com peso $\sigma/10$).
              \end{itemize}
      \end{itemize}
    \end{column}
    \begin{column}{0.44\textwidth}
      \begin{regra}
        \scriptsize
        \tr{\textbf{Estatística final do ajuste:}\\[3pt]
        -- Graus de liberdade: $\mathrm{GL} = \bnum{\num{9702}}$;\\[2pt]
        -- Fator de variância a posteriori: $\hat{\sigma}_0^2 = \bnum{\num{1.509}}$;\\[2pt]
        -- A rede foi particionada em \bnum{8 blocos de Helmert} resolvidos sem nenhuma costura
        residual nas fronteiras regionais. [B]}{\textbf{Final adjustment statistics:}\\[3pt]
        -- Degrees of freedom: $\mathrm{DF} = \bnum{\num{9702}}$;\\[2pt]
        -- A posteriori variance factor: $\hat{\sigma}_0^2 = \bnum{\num{1.509}}$;\\[2pt]
        -- Partitioned into \bnum{8 Helmert blocks} solved without residual boundary seams across
        regions. [B]}
      \end{regra}
    \end{column}
  \end{columns}
\end{frame}

\begin{frame}{\tr{Conexões GPS: o choque entre o antigo e o novo}{GPS Connections: Clashing Old and New}}
  \begin{columns}[onlytextwidth]
    \begin{column}{0.52\textwidth}
      \begin{itemize}
        \setlength{\itemsep}{4pt}
        \scriptsize
        \item \tr{Entre 2004 e 2005, o IBGE realizou \bnum{\num{12} campanhas GPS de campo} para
              ocupar marcos antigos da rede de triangulação. [B]}{Between 2004 and 2005, IBGE carried
              out \bnum{\num{12} field GPS campaigns} reoccupying classical triangulation marks. [B]}
        \item \tr{O resultado causou espanto: quando um marco clássico recebia uma coordenada GPS
              precisa, o ajustamento puxava os vizinhos em até \bnum{\num{9}\,m}, com um raio de
              influência de cerca de \bnum{\num{500}\,km}! [B]}{The result was startling: when an old
              mark received precise GPS coordinates, the adjustment shifted neighbors by up to
              \bnum{\num{9}\,m}, with an influence radius of \bnum{\num{500}\,km}! [B]}
        \item \tr{Isso comprovou que a rede clássica SAD69 estava cheia de deformações elásticas
              locais causadas pelo acúmulo de visadas angulares antigas.}{This proved the classical
              SAD69 network harbored local elastic distortions from accumulated optical angle errors.}
      \end{itemize}
    \end{column}
    \begin{column}{0.44\textwidth}
      \centering
      \begin{tikzpicture}[scale=0.8]
        \draw[thick, corFrame!40] (0,0) circle (2.2);
        \draw[dashed, seta2, corFrame!70] (0,0) -- (1.5,1.6) node[midway, above left, font=\tiny] {$R \approx \num{500}$\,km};
        \fill[mTeal] (0,0) circle (3pt) node[below, font=\scriptsize] {\tr{Vértice GPS}{GPS Vertex}};
        \fill[mRose] (1.1,0.5) circle (2pt) node[right, font=\tiny] {\tr{Marco clássico}{Classical mark}};
        \draw[seta, mRose, thick] (1.1,0.5) -- (0.5,0.7) node[midway, above, font=\tiny] {\tr{até \num{9}\,m}{up to \num{9}\,m}};
      \end{tikzpicture}\\[2pt]
      \tiny \tr{O GPS funcionou não apenas como régua, mas como auditor rigoroso da rede clássica.}{GPS
      acted not merely as a ruler, but as a rigorous auditor of the classical network.}
    \end{column}
  \end{columns}
\end{frame}

\begin{frame}{\tr{O Marco Legal Brasileiro: Decreto 5.334 e R.PR 1/2005}{Brazilian Legal Framework: Decree 5,334 and R.PR 1/2005}}
  \begin{columns}[onlytextwidth]
    \begin{column}{0.48\textwidth}
      \begin{tcolorbox}[colback=mIndigo!8, colframe=mIndigo, title={\tr{Decreto Federal nº 5.334}{Federal Decree no. 5,334}}]
        \scriptsize
        \textbf{\tr{06 de janeiro de 2005 (Presidência da República)}{January 6, 2005 (Presidency of the Republic)}}\\[3pt]
        \tr{Alterou o artigo 21 do Decreto nº 89.817/1984 (Normas Técnicas da Cartografia Nacional):\\[3pt]
        ``Os referenciais planimétrico e altimétrico da cartografia nacional são os que constituem o
        \bnum{Sistema Geodésico Brasileiro}, \emph{definidos pelo IBGE}''. [E]}{Amended Article 21
        of Decree 89,817/1984 (National Cartographic Standards):\\[3pt]
        ``The planimetric and altimetric reference systems of national cartography are those of the
        \bnum{Brazilian Geodetic System}, \emph{as defined by IBGE}''. [E]}
      \end{tcolorbox}
    \end{column}
    \begin{column}{0.48\textwidth}
      \begin{tcolorbox}[colback=mTeal!8, colframe=mTeal, title={\tr{Resolução PR nº 01/2005}{Resolution PR no. 01/2005}}]
        \scriptsize
        \textbf{\tr{25 de fevereiro de 2005 (Presidência do IBGE)}{February 25, 2005 (IBGE Presidency)}}\\[3pt]
        \tr{-- Oficializou o \bnum{SIRGAS2000} como novo referencial oficial do Brasil (GRS80,
        geocêntrico).\\[2pt]
        -- Revogou a R.PR 22/1983 (que fixava o SAD69).\\[2pt]
        -- Estabeleceu um \bnum{período de transição de 10 anos} de convivência com os referenciais
        antigos.}{-- Decreed \bnum{SIRGAS2000} as Brazil's official reference frame (GRS80,
        geocentric).\\[2pt]
        -- Repealed R.PR 22/1983 (which established SAD69).\\[2pt]
        -- Established a \bnum{10-year transition period} allowing dual usage with older datums.}
      \end{tcolorbox}
    \end{column}
  \end{columns}
\end{frame}

\begin{frame}{\tr{25 de fevereiro de 2015: o fim da transição}{February 25, 2015: End of the Transition}}
  \begin{columns}[onlytextwidth]
    \begin{column}{0.54\textwidth}
      \begin{itemize}
        \setlength{\itemsep}{4pt}
        \scriptsize
        \item \tr{Em \bnum{25 de fevereiro de 2015}, encerrou-se rigorosamente o período de transição
              de 10 anos estabelecido pela Resolução PR 01/2005.}{On \bnum{February 25, 2015}, the
              10-year transition period established by Resolution PR 01/2005 officially expired.}
        \item \tr{A partir dessa data:}{From that date forward:}
              \begin{itemize}
                \tiny
                \item \textbf{Córrego Alegre e SAD69} foram definitivamente extintos como referenciais legais;
                \item Todo trabalho cartográfico, cadastral e de engenharia no Brasil deve ser compulsoriamente em \bnum{SIRGAS2000};
                \item O INCRA não aceita mais nenhuma planta de georreferenciamento de imóveis rurais (Lei 10.267) fora do SIRGAS2000.
              \end{itemize}
      \end{itemize}
    \end{column}
    \begin{column}{0.42\textwidth}
      \begin{cuidado}
        \scriptsize
        \tr{\textbf{Atenção jurídica:}\\[3pt]
        Entregar uma planta topográfica, projeto de loteamento ou estudo de desapropriação em SAD69
        hoje é vício técnico grave e causa nulidade documental perante cartórios de registro de imóveis
        e órgãos ambientais!}{\textbf{Legal warning:}\\[3pt]
        Delivering a boundary survey, subdivision plan, or engineering project in SAD69 today is a
        serious technical defect causing legal invalidity in land registries and environmental licensing!}
      \end{cuidado}
    \end{column}
  \end{columns}
\end{frame}

\begin{frame}{\tr{A conta da conversão: os 3 Parâmetros Oficiais}{The Conversion Formula: The 3 Official Parameters}}
  \begin{columns}[onlytextwidth]
    \begin{column}{0.52\textwidth}
      \scriptsize
      \tr{Para converter coordenadas cartesianas de SAD69 para SIRGAS2000, o IBGE determinou em 2005
      e ratificou em 2006 [B, tab. 9] três parâmetros de translação pura geocêntrica:}{To convert
      Cartesian coordinates from SAD69 to SIRGAS2000, IBGE determined in 2005 and reaffirmed in 2006
      [B, tab. 9] three pure geocentric translation parameters:}
      \begin{equation*}
        \begin{pmatrix} X \\ Y \\ Z \end{pmatrix}_{\text{SIRGAS}} =
        \begin{pmatrix} X \\ Y \\ Z \end{pmatrix}_{\text{SAD69}} +
        \begin{pmatrix} \Delta X \\ \Delta Y \\ \Delta Z \end{pmatrix}
      \end{equation*}
      \begin{itemize}
        \setlength{\itemsep}{2pt}
        \tiny
        \item $\Delta X = \bnum{-\num{67.35}}\,\mathrm{m}$;
        \item $\Delta Y = \bnum{+\num{3.88}}\,\mathrm{m}$;
        \item $\Delta Z = \bnum{-\num{38.22}}\,\mathrm{m}$.
      \end{itemize}
      \vspace{0.1cm}
      \tr{Módulo do vetor de deslocamento 3D:}{Magnitude of the 3D shift vector:}
      \begin{equation*}
        |\Delta \vec{X}| = \sqrt{(-67{,}35)^2 + 3{,}88^2 + (-38{,}22)^2} \approx \bnum{\num{77.54}}\,\mathrm{m}
      \end{equation*}
    \end{column}
    \begin{column}{0.44\textwidth}
      \begin{regra}
        \scriptsize
        \tr{\textbf{Por que apenas 3 parâmetros?}\\[3pt]
        O IBGE testou modelos de Helmert de 7 parâmetros (com 3 rotações e escala) em 21, 63 e 113
        estações comuns.\\[3pt]
        Constatou-se que as rotações e a escala eram estatisticamente desprezíveis para o conjunto
        nacional, e um modelo de 7 parâmetros causava instabilidade nos extremos do país. [B]}{
        \textbf{Why only 3 parameters?}\\[3pt]
        IBGE tested 7-parameter Helmert models (3 rotations and scale) across 21, 63, and 113 common
        stations.\\[3pt]
        Rotations and scale were found statistically negligible nationally, while 7 parameters
        induced edge instabilities at border extremes. [B]}
      \end{regra}
    \end{column}
  \end{columns}
\end{frame}

\begin{frame}{\tr{O mito dos ``65 metros'': o deslocamento varia no chão!}{The ``65-Meter'' Myth: Real Shifts Vary by Region!}}
  \begin{columns}[onlytextwidth]
    \begin{column}{0.54\textwidth}
      \scriptsize
      \tr{O vetor 3D de \bnum{\num{77.54}\,m} tem direção fixa no espaço. Ao ser projetado na
      superfície curva da Terra, o deslocamento horizontal varia de \bnum{\num{57}} a \bnum{\num{76}\,m}!
      (calculado por \texttt{sgb\_sirgas2000\_numeros.py}):}{The 3D vector of \bnum{\num{77.54}\,m}
      has a fixed spatial orientation. Projected onto Earth's curved surface, the horizontal shift
      varies from \bnum{\num{57}} to \bnum{\num{76}\,m}! (computed via \texttt{sgb\_sirgas2000\_numeros.py}):}
      \vspace{0.1cm}
      \tiny
      \begin{tabular}{lcccc}
        \toprule
        \textbf{\tr{Local}{City}} & \textbf{$dN$ (m)} & \textbf{$dE$ (m)} & \textbf{$dh$ (m)} & \textbf{\tr{Horiz.}{Horiz.}} \\
        \midrule
        Fortaleza & $-\num{41.65}$ & $-\num{38.92}$ & $-\num{29.50}$ & \bnum{\num{57.00}\,m} \\
        Brasília  & $-\num{49.54}$ & $-\num{47.36}$ & $-\num{12.87}$ & \bnum{\num{68.54}\,m} \\
        Manaus    & $-\num{40.11}$ & $-\num{56.40}$ & $-\num{11.88}$ & \bnum{\num{69.21}\,m} \\
        Curitiba  & $-\num{54.20}$ & $-\num{48.50}$ & $-\num{3.02}$  & \bnum{\num{72.73}\,m} \\
        Chuí      & $-\num{55.22}$ & $-\num{51.80}$ & $+\num{8.06}$  & \bnum{\num{75.72}\,m} \\
        Cruzeiro do Sul & $-\num{40.88}$ & $-\num{63.14}$ & $+\num{4.51}$ & \bnum{\num{75.22}\,m} \\
        \bottomrule
      \end{tabular}
    \end{column}
    \begin{column}{0.42\textwidth}
      \begin{cuidado}
        \scriptsize
        \tr{\textbf{Atenção prática:}\\[3pt]
        Nunca use uma regra fixa de ``subtrair 65 metros''!\\[3pt]
        Em Fortaleza o deslocamento horizontal é de apenas \bnum{\num{57.0}\,m}, enquanto no Chuí e no
        Acre ultrapassa \bnum{\num{75}\,m}.\\[3pt]
        Além disso, a altura elipsoidal ($h$) varia de \bnum{-\num{29.5}\,m} no Nordeste a
        \bnum{+\num{8.1}\,m} no Sul!}{\textbf{Field warning:}\\[3pt]
        Never apply a crude rule like ``subtract 65 meters''!\\[3pt]
        In Fortaleza the horizontal shift is only \bnum{\num{57.0}\,m}, whereas in Chuí and Acre it
        exceeds \bnum{\num{75}\,m}.\\[3pt]
        Furthermore, ellipsoidal height ($h$) shifts from \bnum{-\num{29.5}\,m} in the Northeast to
        \bnum{+\num{8.1}\,m} in the South!}
      \end{cuidado}
    \end{column}
  \end{columns}
\end{frame}

\begin{frame}{\tr{O Método de Grade: ProGriD e o formato NTv2}{Grid Shift Method: ProGriD and NTv2}}
  \begin{columns}[onlytextwidth]
    \begin{column}{0.52\textwidth}
      \begin{itemize}
        \setlength{\itemsep}{4pt}
        \scriptsize
        \item \tr{Os 3 parâmetros tratam apenas da mudança do centro do elipsoide. Mas e as
              \bnum{distorções locais de até \num{13}\,m} da rede clássica?}{The 3 parameters model
              only the shift between ellipsoid centers. But what about the \bnum{local distortions
              up to \num{13}\,m} in the old network?}
        \item \tr{Para compatibilizar cartas topográficas antigas (escalas 1:25.000 e 1:50.000), o
              IBGE desenvolveu o \textbf{ProGriD}.}{To transform legacy map sheets (1:25,000 and
              1:50,000), IBGE developed \textbf{ProGriD}.}
        \item \tr{Baseado no padrão internacional canadense \textbf{NTv2} (\emph{National
              Transformation version 2}):}{Based on Canada's standard \textbf{NTv2} (\emph{National
              Transformation version 2}):}
              \begin{itemize}
                \tiny
                \item Uma grade regular modela os deslocamentos locais $\delta \varphi, \delta \lambda$;
                \item Interpolação bilinear suave na célula da grade;
                \item Modela distorções locais que 3 parâmetros não conseguem absorver.
              \end{itemize}
      \end{itemize}
    \end{column}
    \begin{column}{0.44\textwidth}
      \begin{ideia}
        \scriptsize
        \tr{\textbf{Quando usar cada método?}\\[3pt]
        -- \textbf{3 Parâmetros:} para levantamentos modernos amarrados à rede GPS ou conversões
        globais.\\[3pt]
        -- \textbf{ProGriD (NTv2):} para converter dados históricos digitalizados de cartas antigas
        impressas em papel, onde as distorções da triangulação clássica estão congeladas no desenho.}{
        \textbf{When to use which?}\\[3pt]
        -- \textbf{3 Parameters:} for modern GNSS surveys or general geographic conversions.\\[3pt]
        -- \textbf{ProGriD (NTv2):} for converting digitized legacy maps surveyed from classical
        triangulation, where optical distortions are frozen in the original paper drafting.}
      \end{ideia}
    \end{column}
  \end{columns}
\end{frame}

\begin{frame}{\tr{A crosta não para: a vida útil do SIRGAS2000}{The Crust Never Stops: Lifetime of SIRGAS2000}}
  \begin{columns}[onlytextwidth]
    \begin{column}{0.52\textwidth}
      \begin{itemize}
        \setlength{\itemsep}{4pt}
        \scriptsize
        \item \tr{Lembre-se: o SIRGAS2000 está \textbf{congelado na época 2000,4}.}{Recall: SIRGAS2000
              is \textbf{frozen at epoch 2000.4}.}
        \item \tr{Mas a Placa Sul-Americana desliza a cerca de \bnum{\num{1.2}\,cm/ano} para Noroeste
              em relação ao manto profundo da Terra.}{However, the South American Plate glides at
              about \bnum{\num{1.2}\,cm/yr} northwestward relative to Earth's deep mantle.}
        \item \tr{Hoje (\bnum{2026}), passaram-se \bnum{\num{25.6} anos} desde a campanha de maio de
              2000:}{Today (\bnum{2026}), \bnum{\num{25.6} years} have elapsed since the May 2000
              campaign:}
              \begin{equation*}
                \Delta s \approx \num{1.2}\,\mathrm{cm/ano} \times \num{25.6}\,\mathrm{anos} \approx \bnum{\num{31}\,cm}!
              \end{equation*}
        \item \tr{Se você ligar um receptor GNSS geodésico no campo hoje e medir no ITRF atual, sua
              posição física dista cerca de \bnum{\num{31}\,cm} da coordenada na monografia! [D]}{If
              you turn on a geodetic GNSS receiver in the field today measuring in current ITRF, your
              physical mark sits about \bnum{\num{31}\,cm} away from the datasheet coordinate! [D]}
      \end{itemize}
    \end{column}
    \begin{column}{0.44\textwidth}
      \centering
      \begin{tikzpicture}[scale=0.9]
        \draw[thick, corSist] (0,0) circle (2pt) node[below, font=\tiny] {\tr{Coordenada 2000,4 (no papel)}{Epoch 2000.4 (datasheet)}};
        \draw[seta, corTempo, very thick] (0,0) -- (-1.8,1.2) node[midway, above left, font=\tiny] {$\vec{v} \Delta t \approx \num{31}$\,cm};
        \draw[thick, corFrame] (-1.8,1.2) circle (2pt) node[above, font=\tiny] {\tr{Posição física real em 2026}{Physical position in 2026}};
      \end{tikzpicture}\\[2pt]
      \tiny \tr{O marco viajou \num{31}\,cm sobre a placa tectônica nos últimos 26 anos.}{The physical
      pillar traveled \num{31}\,cm riding the tectonic plate over the past 26 years.}
    \end{column}
  \end{columns}
\end{frame}

\begin{frame}{\tr{Modelos de Velocidade: VEMOS e SIRGAS-CON}{Velocity Models: VEMOS and SIRGAS-CON}}
  \begin{columns}[onlytextwidth]
    \begin{column}{0.54\textwidth}
      \begin{itemize}
        \setlength{\itemsep}{4pt}
        \scriptsize
        \item \tr{Para conciliar a posição medida no GNSS hoje com o referencial oficial congelado em
              2000,4, o SIRGAS calcula e atualiza os modelos de velocidade \textbf{VEMOS} (\emph{Velocity
              Model for SIRGAS}).}{To reconcile modern GNSS observations with the official frame
              frozen at 2000.4, SIRGAS maintains \textbf{VEMOS} (\emph{Velocity Model for SIRGAS}).}
        \item \tr{Gerações de modelos: VEMOS2003, VEMOS2009, VEMOS2015 e \bnum{VEMOS2022}, baseadas em
              mais de \num{400} estações contínuas da rede \textbf{SIRGAS-CON}.}{Model generations:
              VEMOS2003, VEMOS2009, VEMOS2015, and \bnum{VEMOS2022}, derived from over \num{400}
              continuous \textbf{SIRGAS-CON} stations.}
        \item \tr{\textbf{A equação de redução de época:}}{\textbf{Epoch reduction equation:}}
              \begin{equation*}
                \vec{X}(2000{,}4) = \vec{X}(t) - \vec{v} \cdot (t - 2000{,}4)
              \end{equation*}
      \end{itemize}
    \end{column}
    \begin{column}{0.42\textwidth}
      \begin{regra}
        \scriptsize
        \tr{\textbf{Como o IBGE-PPP resolve isso:}\\[3pt]
        Quando você sobe um arquivo RINEX medido ontem no IBGE-PPP, o software calcula primeiro sua
        coordenada na época do rastreio e, em seguida, aplica o modelo de velocidades da placa para
        entregar a coordenada oficial reduzida para \bnum{2000,4}!}{\textbf{How IBGE-PPP handles this:}\\[3pt]
        When you submit a RINEX file recorded yesterday to IBGE-PPP, the server solves your coordinate
        at the survey epoch and then applies crustal velocity models to output coordinates reduced to
        \bnum{2000.4}!}
      \end{regra}
    \end{column}
  \end{columns}
\end{frame}

\begin{frame}{\tr{Debate moderno: manter congelado ou adotar datum dinâmico?}{Modern Debate: Keep Frozen or Adopt Dynamic Datums?}}
  \begin{columns}[onlytextwidth]
    \begin{column}{0.48\textwidth}
      \begin{tcolorbox}[colback=mTeal!8, colframe=mTeal, title={\tr{Modelo Estático (Brasil / SIRGAS)}{Static Frame (Brazil / SIRGAS)}}]
        \scriptsize
        \begin{itemize}
          \setlength{\itemsep}{2pt}
          \item \textbf{\tr{Vantagem:}{Advantage:}} coordenadas de imóveis e mapas permanecem fixas por décadas. Ninguém precisa refazer escrituras a cada ano.
          \item \textbf{\tr{Desafio:}{Challenge:}} com o passar dos anos, a defasagem entre o GNSS de campo e o papel cresce (\num{31}\,cm em 2026, \num{50}\,cm em 2040), exigindo modelos de velocidade.
        \end{itemize}
      \end{tcolorbox}
    \end{column}
    \begin{column}{0.48\textwidth}
      \begin{tcolorbox}[colback=mAmber!8, colframe=mAmber, title={\tr{Modelo Semi-Dinâmico (Austrália)}{Semi-Dynamic Frame (Australia)}}]
        \scriptsize
        \begin{itemize}
          \setlength{\itemsep}{2pt}
          \item \textbf{\tr{Exemplo:}{Example:}} a Austrália desloca-se a $\approx \num{7}$\,cm/ano para Nordeste. Em 2020 abandonou o GDA94 e lançou o \textbf{GDA2020} (\num{1.8}\,m de salto!).
          \item \textbf{\tr{Desafio:}{Challenge:}} exige remapear todo o cadastro territorial nacional e mudar bases GIS a cada atualização de época.
        \end{itemize}
      \end{tcolorbox}
    \end{column}
  \end{columns}
\end{frame}

\begin{frame}{\tr{Sua vez: reduzindo uma coordenada para a época oficial}{Your Turn: Reducing a Coordinate to Official Epoch}}
  \begin{columns}[T]
    \begin{column}{0.50\textwidth}
      \begin{ideia}[title={\tr{Problema}{Problem}}]
        \scriptsize
        \tr{Um receptor geodésico mediu em \bnum{2025,4} a coordenada $X$ de um marco no ITRF como:
        $X(\num{2025.4}) = \num{2754120.350}$\,m.
        \begin{enumerate}
          \setlength{\itemsep}{2pt}
          \item[(a)] Se a velocidade da placa nessa componente for $v_X = -\num{0.010}$\,m/ano, qual é
                o deslocamento acumulado em \num{25} anos?
          \item[(b)] Qual é a coordenada oficial $X$ que deve ser impressa na monografia do SIRGAS2000
                (época 2000,4)?
        \end{enumerate}}{A geodetic receiver surveyed in \bnum{2025.4} the $X$ coordinate of a mark in
        ITRF as: $X(\num{2025.4}) = \num{2754120.350}$\,m.
        \begin{enumerate}
          \setlength{\itemsep}{2pt}
          \item[(a)] If plate velocity is $v_X = -\num{0.010}$\,m/yr, what is the shift over \num{25} years?
          \item[(b)] What is the official $X$ coordinate for the SIRGAS2000 datasheet (epoch 2000.4)?
        \end{enumerate}}
      \end{ideia}
    \end{column}
    \begin{column}{0.46\textwidth}
      \begin{regra}[title={\tr{Resolução passo a passo}{Step-by-step Solution}}]
        \scriptsize\raggedright
        (a) $\Delta X = v_X \cdot (t - t_0)$\\[2pt]
        $\phantom{(a)} = (-\num{0.010}\,\mathrm{m/ano}) \times (\num{25}\,\mathrm{anos}) = \bnum{-\num{0.250}\,\mathrm{m}}$.
        \medskip

        (b) $X(2000{,}4) = X(t) - \Delta X$\\[2pt]
        $\phantom{(b)} = \num{2754120.350} - (-\num{0.250})$\\[2pt]
        $\phantom{(b)} = \bnum{\num{2754120.600}}\,\mathrm{m}$.
        \medskip

        \tr{\textbf{Conclusão:} Uma diferença de \bnum{\num{25}\,cm} entre a época do rastreio e a
        coordenada legal oficial!}{\textbf{Conclusion:} A difference of \bnum{\num{25}\,cm} between
        survey epoch and official legal coordinate!}
      \end{regra}
    \end{column}
  \end{columns}
\end{frame}

\begin{frame}{\tr{A grande ponte para a Parte 2: a água não corre pelo elipsoide!}{The Grand Bridge to Part 2: Water Doesn't Flow on the Ellipsoid!}}
  \begin{columns}[onlytextwidth]
    \begin{column}{0.52\textwidth}
      \begin{itemize}
        \setlength{\itemsep}{4pt}
        \scriptsize
        \item \tr{O SIRGAS2000 resolveu com perfeição a geometria tridimensional do país:
              \bnum{latitude $\varphi$, longitude $\lambda$ e altura elipsoidal $h$}.}{SIRGAS2000
              flawlessly solved Brazil's 3D geometry: \bnum{latitude $\varphi$, longitude $\lambda$,
              and ellipsoidal height $h$}.}
        \item \tr{Mas na engenharia, saneamento, canais e barragens, \alert{a água não se importa com o
              elipsoide matemático}!}{However, in civil engineering, sanitation, canals, and dams,
              \alert{water does not follow the mathematical ellipsoid}!}
        \item \tr{A água escorre pelas superfícies equipotenciais da \textbf{gravidade terrestre}:}{
              Water flows along equipotential surfaces of \textbf{Earth's gravity field}:}
              \begin{itemize}
                \tiny
                \item Dois pontos com a mesma altura elipsoidal ($h_A = h_B$) podem ter altitudes físicas diferentes e a água escorrer de um para o outro!
              \end{itemize}
        \item \tr{A altitude da vida real ($H$) vem da \textbf{Rede Altimétrica Fundamental} amarrada
              ao Marégrafo de Imbituba (SC).}{Physical height ($H$) comes from the \textbf{Vertical
              Leveling Network} tied to Imbituba Tide Gauge (SC).}
      \end{itemize}
    \end{column}
    \begin{column}{0.44\textwidth}
      \centering
      \begin{tikzpicture}[scale=0.9]
        \draw[thick, corSist] (-2,0) arc[start angle=140, end angle=40, radius=3.0] node[right, font=\tiny] {\tr{Elipsoide GRS80}{GRS80 Ellipsoid}};
        \draw[thick, corMar] (-2,0.6) arc[start angle=145, end angle=35, radius=2.7] node[right, font=\tiny] {\tr{Geoide / Mar Médio}{Geoid / Mean Sea Level}};
        \draw[seta2, corErro, thick] (0,0.52) -- (0,1.25) node[midway, right, font=\tiny] {$N$ \tr{(ondulação)}{undulation}};
        \draw[seta, thick, mTeal] (-0.8,1.6) -- (0.8,1.4) node[midway, above, font=\tiny] {\tr{Fluxo d'água}{Water flow}};
      \end{tikzpicture}\\[2pt]
      \tiny \tr{A ponte entre $h$ (SIRGAS) e $H$ (Imbituba) é o modelo de geoide.}{The bridge between
      $h$ (SIRGAS) and $H$ (Imbituba) is the geoid model.}
    \end{column}
  \end{columns}
\end{frame}

\begin{frame}{\tr{Resumo da Parte 1: As Quatro Aulas num Olhar}{Part 1 Summary: The Four Classes at a Glance}}
  \centering
  \scriptsize
  \begin{tabular}{cp{10.5cm}}
    \toprule
    \textbf{\tr{Aula}{Class}} & \textbf{\tr{A Ideia Central da Aula}{Core Idea of the Class}} \\
    \midrule
    \textbf{1} & \tr{\textbf{A Herança:} Os datums clássicos (Córrego Alegre, SAD69) amarravam o elipsoide ao relevo mineiro (Chuá). Sem geocentro, a rede clássica acumulou até \bnum{\num{13}\,m} de distorção nos extremos.}{\textbf{The Heritage:} Classical datums (Córrego Alegre, SAD69) pinned ellipsoids to Minas Gerais (Chuá). Without geocenter, classical triangulation piled up \bnum{\num{13}\,m} distortion at border extremes.} \\[3pt]
    \textbf{2} & \tr{\textbf{A Idealização:} Em Assunção (1993), decidiu-se pregar o referencial no centro de massa da Terra (ITRS, GRS80, GPS). O SIRGAS95 provou a precisão milimétrica em 58 estações continentais.}{\textbf{The Design:} In Asunción (1993), geodesists chose to anchor the frame to Earth's center of mass (ITRS, GRS80, GPS). SIRGAS95 demonstrated millimeter accuracy across 58 continental sites.} \\[3pt]
    \textbf{3} & \tr{\textbf{A Realização:} Em maio de 2000, 184 estações (21 no Brasil) rastrearam simultaneamente. O processamento independente em três centros e o ajustamento com injunções geraram o SIRGAS2000 na época congelada \bnum{2000,4}.}{\textbf{The Realization:} In May 2000, 184 stations (21 in Brazil) tracked simultaneously. Processing in three independent centers and constrained adjustment yielded SIRGAS2000 frozen at epoch \bnum{2000.4}.} \\[3pt]
    \textbf{4} & \tr{\textbf{A Adoção e Vida Útil:} Decreto 5.334 e R.PR 1/2005 (transição até 2015). O reajuste de 2006 (6265 estações). A placa deriva \bnum{\num{1.2}\,cm/ano} (\num{31}\,cm em 2026), corrigida por modelos VEMOS.}{\textbf{Adoption and Lifetime:} Decree 5,334 and R.PR 1/2005 (transition ending 2015). The 2006 re-adjustment (6,265 stations). Plate drifts \bnum{\num{1.2}\,cm/yr} (\num{31}\,cm in 2026), modeled by VEMOS.} \\
    \bottomrule
  \end{tabular}
  \vspace{0.25cm}
  \begin{ideia}
    \centering \footnotesize
    \tr{\textbf{Continua na Parte 2:} O Sistema Geodésico Brasileiro --- A Rede Altimétrica, o Nivelamento e o Geoide!}{\textbf{Continues in Part 2:} The Brazilian Geodetic System --- The Leveling Network, Heights, and the Geoid!}
  \end{ideia}
\end{frame}

% ==================================================================
\appendix

\begin{frame}[standout]
  \tr{Apêndice Técnico}{Technical Appendix}\\[6pt]
  {\normalsize \tr{Formulações matemáticas, linha do tempo e fontes primárias}{Mathematical formulations, timeline and primary sources}}
\end{frame}

\begin{frame}{\tr{A1 · Definição do ITRS e do SIRGAS}{A1 · ITRS and SIRGAS Definition}}
  \scriptsize
  \begin{itemize}
    \setlength{\itemsep}{3pt}
    \item \textbf{\tr{Origem:}{Origin:}} Centro de massa da Terra total (incluindo oceanos e atmosfera).
    \item \textbf{\tr{Escala:}{Scale:}} Determinada no sistema próprio de coordenadas relativísticas terrestres (TT / TCG).
    \item \textbf{\tr{Orientação:}{Orientation:}} Polo convencional terrestre (CIP) e meridiano de referência internacional (IRM), alinhados às convenções do IERS / BIH 1984.0.
    \item \textbf{\tr{Evolução temporal (Condição NNR):}{Time evolution (NNR Condition):}}
          A rotação líquida residual da crosta terrestre em relação ao manto profundo deve ser nula:
          \begin{equation*}
            \int_{\text{Terra}} (\vec{r} \times \vec{v}) \, dm = \vec{0}
          \end{equation*}
    \item \textbf{\tr{Elipsoide associado:}{Associated Ellipsoid:}} GRS80 (\emph{Geodetic Reference System 1980}):
          \begin{equation*}
            a = \bnum{\num{6378137}}\,\mathrm{m}, \quad 1/f = \bnum{\num{298.257222101}}, \quad GM = \num{3986005}\times10^8\,\mathrm{m^3/s^2}
          \end{equation*}
  \end{itemize}
\end{frame}

\begin{frame}{\tr{A2 · Ajustamento por Mínimos Quadrados com Injunções}{A2 · Least Squares Adjustment with Constraints}}
  \scriptsize
  \tr{Seja o modelo de observações linearizado $v = A \hat{x} - l$ com pesos $P = \sigma_0^2 \Sigma^{-1}$, e as equações de injunção nas coordenadas conhecidas $v_C = C \hat{x} - w$ com pesos $P_C = \sigma_0^2 \Sigma_C^{-1}$.}{
  Given the linearized observation model $v = A \hat{x} - l$ with weights $P = \sigma_0^2 \Sigma^{-1}$, and constraint equations on known coordinates $v_C = C \hat{x} - w$ with weights $P_C = \sigma_0^2 \Sigma_C^{-1}$.}
  \vspace{0.15cm}
  \begin{columns}[onlytextwidth]
    \begin{column}{0.48\textwidth}
      \begin{tcolorbox}[colback=mStoneLight, colframe=mStone, title={\tr{Sistema Normal Combinado}{Combined Normal System}}]
        \scriptsize
        \begin{equation*}
          (A^T P A + C^T P_C C) \hat{x} = A^T P l + C^T P_C w
        \end{equation*}
        \tr{Definindo $N = A^T P A$ e $u = A^T P l$:}{Defining $N = A^T P A$ and $u = A^T P l$:}
        \begin{equation*}
          \hat{x} = (N + C^T P_C C)^{-1} (u + C^T P_C w)
        \end{equation*}
      \end{tcolorbox}
    \end{column}
    \begin{column}{0.48\textwidth}
      \begin{tcolorbox}[colback=mTeal!8, colframe=mTeal, title={\tr{Variância a Posteriori}{A Posteriori Variance}}]
        \scriptsize
        \begin{equation*}
          \hat{\sigma}_0^2 = \frac{v^T P v + v_C^T P_C v_C}{(n + n_C) - u}
        \end{equation*}
        \tr{Onde:}{Where:}
        \begin{itemize}
          \tiny
          \item $n$: número de observações de campo;
          \item $n_C$: número de equações de injunção;
          \item $u$: número de parâmetros incógnitos estimados.
        \end{itemize}
      \end{tcolorbox}
    \end{column}
  \end{columns}
\end{frame}

\begin{frame}{\tr{A3 · Fórmulas de Transformação SAD69 $\to$ SIRGAS2000}{A3 · Transformation Formulas SAD69 $\to$ SIRGAS2000}}
  \scriptsize
  \tr{A transformação oficial brasileira recomendada pelo IBGE opera no espaço geocêntrico 3D:}{The
  official Brazilian transformation recommended by IBGE operates in 3D geocentric space:}
  \begin{equation*}
    \begin{pmatrix} X \\ Y \\ Z \end{pmatrix}_{\text{SIRGAS2000}} =
    \begin{pmatrix} X \\ Y \\ Z \end{pmatrix}_{\text{SAD69}} +
    \begin{pmatrix} \Delta X \\ \Delta Y \\ \Delta Z \end{pmatrix}, \quad
    \begin{cases}
      \Delta X = -\num{67.35}\,\mathrm{m} \\
      \Delta Y = +\num{3.88}\,\mathrm{m} \\
      \Delta Z = -\num{38.22}\,\mathrm{m}
    \end{cases}
  \end{equation*}
  \vspace{0.1cm}
  \tr{\textbf{Relação com os parâmetros elipsoidais:}}{ \textbf{Relationship with ellipsoidal parameters:}}
  \begin{itemize}
    \tiny
    \item SAD69 (UGGI67): $a_1 = \num{6378160}$\,m, $1/f_1 = \num{298.25}$;
    \item SIRGAS2000 (GRS80): $a_2 = \num{6378137}$\,m ($\Delta a = a_2 - a_1 = -\num{23}$\,m), $1/f_2 = \num{298.257222101}$ ($\Delta f = -\num{0.000000081}$).
    \item Para conversão analítica direta em latitude e longitude geodésicas sem passar por ECEF, usa-se a fórmula de \textbf{Molodensky Abreviada} (implementada no simulador \texttt{sistemas\_coordenadas/}).
  \end{itemize}
\end{frame}

\begin{frame}{\tr{A4 · Propagação Temporal com Velocidades de Placa}{A4 · Epoch Propagation with Plate Velocities}}
  \scriptsize
  \tr{A posição tridimensional de uma estação no instante $t$ é modelada por cinemática linear:}{The
  3D position of a station at epoch $t$ is modeled via linear kinematics:}
  \begin{equation*}
    \vec{X}(t) = \vec{X}(t_0) + \vec{v} \cdot (t - t_0) + \sum_{k} \Delta \vec{X}_k \, H(t - t_k) + \vec{\varepsilon}(t)
  \end{equation*}
  \begin{columns}[onlytextwidth]
    \begin{column}{0.48\textwidth}
      \begin{itemize}
        \tiny
        \item $\vec{X}(t_0)$: posição de referência na época oficial ($t_0 = 2000{,}4$);
        \item $\vec{v} = (v_X, v_Y, v_Z)^T$: velocidade linear constante da placa (fornecida pelos modelos VEMOS);
        \item $\Delta \vec{X}_k$: saltos co-sísmicos ou mudanças de antena no tempo $t_k$ ($H$: função degrau Heaviside);
        \item $\vec{\varepsilon}(t)$: variações sazonais periódicas (carga hidrológica, marés oceânicas).
      \end{itemize}
    \end{column}
    \begin{column}{0.48\textwidth}
      \begin{tcolorbox}[colback=mIndigo!8, colframe=mIndigo, title={\tr{Para o usuário prático}{For practical users}}]
        \scriptsize \tr{Para levar uma medição GNSS de hoje ($t$) para a época oficial $t_0 = 2000{,}4$:\\[3pt]
        \centering
        $\vec{X}(2000{,}4) = \vec{X}(t) - \vec{v} \cdot (t - 2000{,}4)$}{To reduce today's survey ($t$)
        to official epoch $t_0 = 2000.4$:\\[3pt]
        \centering
        $\vec{X}(2000.4) = \vec{X}(t) - \vec{v} \cdot (t - 2000.4)$}
      \end{tcolorbox}
    \end{column}
  \end{columns}
\end{frame}

\begin{frame}{\tr{A5 · Linha do Tempo da Geodésia Brasileira (1924--2026)}{A5 · Brazilian Geodesy Timeline (1924--2026)}}
  \tiny
  \begin{tabular}{llp{9.5cm}}
    \toprule
    \textbf{\tr{Ano}{Year}} & \textbf{\tr{Marco}{Milestone}} & \textbf{\tr{Significado Técnico}{Technical Significance}} \\
    \midrule
    \textbf{1924} & Hayford 1924 & Adotado no datum Córrego Alegre ($a=\num{6378388}$\,m, $1/f=\num{297}$, $\xi=\eta=0$). \\
    \textbf{1969} & SAD69 & UGGI67 ($a=\num{6378160}$\,m, $1/f=\num{298.25}$), datum no Vértice Chuá (MG). \\
    \textbf{1991} & Chegada do GPS & IBGE adquire os primeiros \num{4} receptores geodésicos GPS de dupla frequência. \\
    \textbf{1993} & Assunção & Conferência cria o SIRGAS (IAG, IPGH, DMA) com os Grupos de Trabalho I e II. \\
    \textbf{1995} & SIRGAS95 & Campanha com \num{58} estações continentais (\num{11} no Brasil). ITRF94 época 1995,4. \\
    \textbf{1996} & Ajuste SAD69 / RBMC & Ajuste global clássico (4760 estações); inauguração da RBMC (Curitiba e Pres. Prudente). \\
    \textbf{1997} & GT III (Rio de Janeiro) & IAG cria o Grupo de Trabalho III (Datum Vertical e Altimetria). \\
    \textbf{2000} & SIRGAS2000 / PMRG & Campanha de maio com \num{184} estações; época \bnum{2000,4}; criação formal do PMRG. \\
    \textbf{2005} & Decreto 5.334 / R.PR 1 & Oficialização legal do SIRGAS2000 no Brasil com período de transição de 10 anos. \\
    \textbf{2006} & Ajuste da Rede Planimétrica & 6.265 estações ajustadas simultaneamente com \num{21} injunções SIRGAS ($\sigma/10$). \\
    \textbf{2015} & Fim da Transição & 25/02/2015: encerramento definitivo do SAD69. Uso exclusivo de SIRGAS2000. \\
    \textbf{2026} & Hoje & Crosta acumulou \num{31}\,cm de deriva; uso rotineiro de IBGE-PPP e modelos VEMOS2022. \\
    \bottomrule
  \end{tabular}
\end{frame}

\begin{frame}{\tr{A6 · Glossário Conceitual Bilíngue}{A6 · Bilingual Conceptual Glossary}}
  \tiny
  \begin{columns}[onlytextwidth]
    \begin{column}{0.48\textwidth}
      \begin{itemize}
        \setlength{\itemsep}{2pt}
        \item \textbf{Datum Topocêntrico / Topocentric Datum:} referencial cujo elipsoide tangencia o geoide num ponto da superfície (Córrego Alegre, Chuá); centro não coincide com o geocentro.
        \item \textbf{Datum Geocêntrico / Geocentric Datum:} referencial com origem no centro de massa da Terra (ITRS, SIRGAS2000, WGS 84).
        \item \textbf{Sistema de Referência / Reference System:} definição teórica e física do espaço (regras, constantes $a, f, GM$, condição NNR).
        \item \textbf{Estrutura de Referência (Frame) / Reference Frame:} materialização física e numérica do sistema (lista de coordenadas e velocidades dos marcos em determinada época).
      \end{itemize}
    \end{column}
    \begin{column}{0.48\textwidth}
      \begin{itemize}
        \setlength{\itemsep}{2pt}
        \item \textbf{Rede Passiva / Passive Network:} marcos físicos cravados no solo (pilares, chapas de bronze) sem rastreamento contínuo.
        \item \textbf{Rede Ativa / Active Network:} estações com receptores GNSS permanentes transmitindo dados ininterruptos (RBMC).
        \item \textbf{Injunção / Constraint:} equação que fixa ou restringe com determinado peso as coordenadas de estações conhecidas no ajustamento.
        \item \textbf{Época de Referência / Reference Epoch:} data específica para a qual as coordenadas de um frame estático são calculadas e congeladas (ex.: 2000,4).
      \end{itemize}
    \end{column}
  \end{columns}
\end{frame}

\begin{frame}{\tr{A7 · Siglas e Instituições Fundamentais}{A7 · Key Acronyms and Institutions}}
  \tiny
  \begin{tabular}{lp{11cm}}
    \toprule
    \textbf{\tr{Sigla}{Acronym}} & \textbf{\tr{Nome e Significado}{Name and Meaning}} \\
    \midrule
    \textbf{SGB} & \tr{Sistema Geodésico Brasileiro (mantido e regulamentado pelo IBGE).}{Brazilian Geodetic System (maintained and regulated by IBGE).} \\
    \textbf{SIRGAS} & \tr{Sistema de Referência Geodésico para as Américas (órgão intergovernamental científico da IAG e IPGH).}{Geodetic Reference System for the Americas.} \\
    \textbf{ITRS / ITRF} & \emph{International Terrestrial Reference System / Frame} (IERS, padrão científico mundial). \\
    \textbf{IERS} & \emph{International Earth Rotation and Reference Systems Service} (Frankfurt / Paris). \\
    \textbf{IAG} & \emph{International Association of Geodesy} (\tr{Associação Internacional de Geodésia}{International Association of Geodesy}). \\
    \textbf{IPGH} & \tr{Instituto Pan-Americano de Geografia e História (órgão especializado da OEA).}{Pan-American Institute of Geography and History (OAS specialized agency).} \\
    \textbf{IGS} & \emph{International GNSS Service} (\tr{serviço global de órbitas e relógios de satélites}{global satellite orbit and clock service}). \\
    \textbf{DGFI-TUM} & \emph{Deutsches Geodätisches Forschungsinstitut} (\tr{Universidade Técnica de Munique, Alemanha}{Technical University of Munich, Germany}). \\
    \textbf{DMA / NGA} & \emph{Defense Mapping Agency} (\tr{EUA, posteriormente NIMA e atual NGA}{USA, later NIMA and now NGA}). \\
    \textbf{IBGE} & \tr{Instituto Brasileiro de Geografia e Estatística (autoridade geodésica e cartográfica do Brasil).}{Brazilian Institute of Geography and Statistics.} \\
    \textbf{RBMC} & \tr{Rede Brasileira de Monitoramento Contínuo dos Sistemas GNSS (rede ativa oficial do IBGE).}{Brazilian Continuous Monitoring Network of GNSS Systems.} \\
    \textbf{SINEX} & \emph{Solution INdependent EXchange} (\tr{formato neutro internacional para intercâmbio de equações normais geodésicas}{neutral international format for exchanging geodetic normal equations}). \\
    \bottomrule
  \end{tabular}
\end{frame}

\begin{frame}{\tr{A8 · Fontes Primárias e Conferência de Fatos}{A8 · Primary Sources and Fact Verification}}
  \tiny
  \begin{itemize}
    \setlength{\itemsep}{2pt}
    \item \textbf{[A]} IBGE (\num{2000}), \emph{SIRGAS Final Report, Working Groups I and II}. Rio de Janeiro: IBGE / DGFI / NIMA. (Campanhas de 1995 e 2000, 58 e 184 estações, ITRF94 e ITRF2000).
    \item \textbf{[B]} IBGE (\num{2006}), \emph{\tr{Ajustamento da Rede Planimétrica Brasileira em SIRGAS2000}{Adjustment of the Brazilian Planimetric Network in SIRGAS2000}}. Rio de Janeiro: Diretoria de Geociências, Coordenação de Geodésia. (6.265 estações, 22.232 equações, parâmetros de translação da tab.\ 9, peso $\sigma/10$).
    \item \textbf{[C]} IBGE (\num{1996}/\num{2002}), \emph{\tr{Ajustamento da Rede Planimétrica do SGB}{Adjustment of the SGB Planimetric Network}}. Rio de Janeiro: IBGE. (Córrego Alegre, SAD69, histórico de triangulação, distorção de até \num{13}\,m no BDG).
    \item \textbf{[D]} SIRGAS (\num{2026}), \texttt{sirgas.ipgh.org}: Seções \emph{About us}, \emph{SIRGAS2000}, \emph{Working Groups I, II, III}, \emph{SIRGAS-CON}, \emph{VEMOS Models}.
    \item \textbf{[E]} \tr{Presidência da República (\num{2005}), \emph{Decreto nº 5.334, de 6 de janeiro de 2005}. Altera o art.\ 21 do Decreto nº 89.817/1984.}{Presidency of the Republic of Brazil (\num{2005}), \emph{Decree no.\ 5,334 of January 6, 2005}.}
    \item \textbf{[F]} \tr{IBGE (\num{2005}), \emph{Resolução do Presidente do IBGE nº 1/2005, de 25 de fevereiro de 2005}. Oficialização do SIRGAS2000 no Brasil e transição de 10 anos.}{IBGE (\num{2005}), \emph{Resolution PR no.\ 1/2005 of February 25, 2005}.}
    \item \tr{Valores numéricos e parâmetros recalculados neste repositório: \texttt{sgb\_sirgas2000\_numeros.py}.}{Derived numbers and parameters recomputed in this repository: \texttt{sgb\_sirgas2000\_numeros.py}.}
  \end{itemize}
  \medskip
  \normalsize
  \centering
  \tr{Os fatos, números e legislações desta aula foram conferidos na íntegra em 13/10/2026.}{All facts, numbers, and legislation in this lecture were verified in full on 13/10/2026.}
\end{frame}

\end{document}
